\documentclass[11pt,a4paper]{article}

%% ===== Core Packages =====
\usepackage{fontspec}
\usepackage{polyglossia}
\setmainlanguage{english}
\usepackage{xeCJK}
\usepackage{amsmath,amssymb,amsthm}
\usepackage{geometry}
\usepackage{graphicx}
\usepackage{xcolor}
\usepackage{titlesec}
\usepackage{fancyhdr}
\usepackage{enumitem}
\usepackage{array,booktabs,longtable,multirow}
\usepackage{hyperref}
\usepackage{tocloft}
\usepackage{indentfirst}
\usepackage{tikz}
\usepackage{paracol}

%% ===== Page Geometry (wider for parallel columns) =====
\geometry{
  paperwidth=210mm,paperheight=297mm,
  left=18mm,right=18mm,top=25mm,bottom=25mm,
  headheight=14pt,footskip=30pt
}

%% ===== Font Configuration =====
\setmainfont{Linux Libertine O}
\setsansfont{Linux Biolinum O}
\setmonofont{DejaVu Sans Mono}

%% ===== CJK Font Setup (as specified) =====
\setCJKmainfont[AutoFakeBold=true,AutoFakeSlant=true]{Noto Sans CJK SC}

%% ===== Colors =====
\definecolor{titlecolor}{RGB}{45,55,72}
\definecolor{sectioncolor}{RGB}{55,65,81}
\definecolor{notecolor}{RGB}{120,53,15}
\definecolor{rulecolor}{RGB}{180,160,140}
\definecolor{coverblue}{RGB}{35,55,85}
\definecolor{chinesebg}{RGB}{250,248,245}
\definecolor{englishbg}{RGB}{245,248,250}
\definecolor{vlinecolor}{RGB}{200,195,190}

%% ===== Column Ratio =====
\columnratio{0.48,0.48}
\setlength{\columnsep}{0.04\textwidth}

%% ===== Section Formatting =====
\titleformat{\section}
  {\normalfont\Large\bfseries\color{sectioncolor}}%
  {\thesection}{1em}{}
\titleformat{\subsection}
  {\normalfont\large\bfseries\color{sectioncolor}}
  {\thesubsection}{1em}{}
\titleformat{\subsubsection}
  {\normalfont\normalsize\bfseries\color{sectioncolor}}
  {\thesubsubsection}{1em}{}

%% ===== Header/Footer =====
\pagestyle{fancy}
\fancyhf{}
\fancyhead[L]{\small\textit{Bilingual Edition --- Ceyuan Haijing Fenlei Shishu}}
\fancyhead[R]{\small\thepage}
\renewcommand{\headrulewidth}{0.4pt}
\renewcommand{\footrulewidth}{0pt}

%% ===== Custom Commands =====
\newcommand{\longline}{\noindent\rule{\textwidth}{0.4pt}}
\newcommand{\shortline}{\noindent\rule{0.5\textwidth}{0.4pt}\par}

%% Chinese structural markers (left column)
\newcommand{\wen}[1]{\noindent\textbf{問} #1\par\smallskip}
\newcommand{\da}[1]{\noindent\textbf{答曰} #1\par\smallskip}
\newcommand{\shu}[1]{\noindent\textbf{術曰} #1\par\smallskip}
\newcommand{\shi}[1]{\noindent\textbf{釋曰} #1\par\smallskip}

%% English structural markers (right column)
\newcommand{\wenE}[1]{\noindent\textbf{Q.} #1\par\smallskip}
\newcommand{\daE}[1]{\noindent\textbf{A.} #1\par\smallskip}
\newcommand{\shuE}[1]{\noindent\textbf{Method:} #1\par\smallskip}
\newcommand{\shiE}[1]{\noindent\textbf{Explanation:} #1\par\smallskip}

%% Bilingual section headers
\newcommand{\bilingsec}[2]{%
  \switchcolumn*\noindent
  \begin{minipage}[t]{0.48\textwidth}\raggedright\bfseries\color{sectioncolor}#1\end{minipage}%
  \hfill%
  \begin{minipage}[t]{0.48\textwidth}\raggedright\bfseries\color{sectioncolor}#2\end{minipage}%
  \par\medskip}

%% Bilingual subsection headers
\newcommand{\bilingsubsec}[2]{%
  \switchcolumn*\noindent
  \begin{minipage}[t]{0.48\textwidth}\raggedright\bfseries #1\end{minipage}%
  \hfill%
  \begin{minipage}[t]{0.48\textwidth}\raggedright\bfseries #2\end{minipage}%
  \par\medskip}

%% Problem environment (bilingual)
\newcounter{problem}[section]
\newenvironment{bilingproblem}[2][]{%
  \refstepcounter{problem}%
  \switchcolumn*\noindent
  \begin{minipage}[t]{0.48\textwidth}\textbf{第\theproblem 題#1}\quad\small #2\end{minipage}%
  \hfill%
  \begin{minipage}[t]{0.48\textwidth}\textbf{Problem \theproblem}\quad\small #2\end{minipage}%
  \par\medskip
  \nopagebreak
}{\medskip\switchcolumn*\noindent\rule{0.48\textwidth}{0.2pt}\hfill\rule{0.48\textwidth}{0.2pt}\par\medskip}

%% TOC Depth
\setcounter{tocdepth}{3}
\setcounter{secnumdepth}{3}

%% Hyperref Setup
\hypersetup{
  colorlinks=true,
  linkcolor=sectioncolor,
  citecolor=sectioncolor,
  urlcolor=sectioncolor,
  pdfencoding=auto
}

%% Prevent overfull hbox warnings in narrow columns
\tolerance=1000
\emergencystretch=3em
\hyphenpenalty=100

\newcommand{\vs}{\vspace{0.5em}}
\begin{document}

%% ========================================================================
%% TITLE PAGE
%% ========================================================================
\thispagestyle{empty}
\begin{center}
\vspace*{1.5cm}
{\fontsize{32}{40}\selectfont
\textcolor{titlecolor}{測圓海鏡分類釋術}}

\vspace{0.8cm}
{\Large\textcolor{sectioncolor}{Ceyuan Haijing Fenlei Shishu}}

\vspace{0.3cm}
{\normalsize\textcolor{sectioncolor}{(Classified Methods of the Sea Mirror of Circle Measurements)}}

\vspace{1cm}
{\Large\textcolor{sectioncolor}{Volumes I--III \quad 卷一至卷三}}

\vspace{1.5cm}
\shortline

\vspace{0.8cm}
{\large
\begin{tabular}{rl}
\textbf{Original Author:} & (元) 李冶 撰 \\[-0.2em]
                          & Li Ye (Yuan Dynasty, 1192--1279) \\[0.5em]
\textbf{Compiler:}        & (明) 顧應祥 分類釋術 \\[-0.2em]
                          & Gu Yingxiang (Ming Dynasty, 1483--1565) \\
\end{tabular}
}

\vspace{1cm}
{\normalsize\textcolor{notecolor}{%
Bilingual English--Chinese Edition\\[0.3em]
Original Chinese with Parallel English Translation\\[0.3em]
Preserving the 問 / 答 / 術 / 釋 Structure
}}

\vspace{1.5cm}
\shortline

\vspace{0.8cm}
{\normalsize\textcolor{sectioncolor}{%
From the Siku Quanshu Huiyao (四庫全書薈要) Collection\\[0.3em]
Zi Bu (子部) --- Mathematical and Scientific Works}}

\vspace{0.8cm}
{\small\textcolor{sectioncolor}{Modern Bilingual Typeset Edition}}

\end{center}
\clearpage

%% ========================================================================
%% COLOPHON PAGE
%% ========================================================================
\thispagestyle{empty}
\begin{center}
\vspace*{2.5cm}
{\LARGE 欽定四庫全書薈要}

\vspace{0.3cm}
{\normalsize\textit{Imperial Collection: Essentials of the Complete Library of the Four Treasuries}}

\vspace{1.2cm}
{\Large 子部}

\vspace{0.3cm}
{\normalsize\textit{Philosophy and Science Division}}

\vspace{1.5cm}
{\Large 測圓海鏡分類釋術卷一至三}

\vspace{0.3cm}
{\normalsize\textit{Ceyuan Haijing Fenlei Shishu, Volumes 1--3}}

\vspace{1.5cm}
\shortline

\vspace{0.8cm}
{\normalsize
\begin{tabular}{ll}
詳校官主事 & 臣陳永 \\
\textit{Chief Examiner:} & \textit{Chen Yong} \\
\end{tabular}
}

\vspace{2cm}
{\small\textcolor{notecolor}{%
This edition was prepared for the Qianlong Emperor's imperial library,\\
the Siku Quanshu Huiyao (Essentials of the Complete Library of the Four Treasuries).
}}

\end{center}
\clearpage

%% ========================================================================
%% TABLE OF CONTENTS
%% ========================================================================
\tableofcontents
\clearpage

%% ========================================================================
%% BILINGUAL LAYOUT BEGINS: Chinese (left) / English (right)
%% ========================================================================

%% ========================================================================
%% IMPERIAL SUMMARY (提要)
%% ========================================================================
\section*{提要 / Imperial Summary}
\addcontentsline{toc}{section}{提要 / Imperial Summary}
\markboth{提要 / Imperial Summary}{提要 / Imperial Summary}

\begin{paracol}{2}

\noindent\textbf{提要}

\medskip

臣等謹案：\textbf{測圓海鏡分類釋術}十卷，元翰林學士\textbf{李冶}撰，明刑部尚書\textbf{顧應祥}分類釋術。

\switchcolumn

\noindent\textbf{Imperial Summary}

\medskip

Your servants respectfully report: The \textit{Ceyuan Haijing Fenlei Shishu} (Classified Methods of the Sea Mirror of Circle Measurements), in ten volumes, was composed by \textbf{Li Ye}, Hanlin Academician of the Yuan Dynasty, and classified with explanations by \textbf{Gu Yingxiang}, Minister of Justice of the Ming Dynasty.

\switchcolumn

冶書前列\textbf{勾股總圖}，以天地日月山川旦夕朱青泛泉等字及四方五行八卦之名於圖線之交，詳為標識。又於交處所成各形，以通商功之類。是以形求積則方田商功之類是也，以積求形則少廣勾股之類是也。以形求積者，先得其形而復求其積，故其為術也易。以積求形則先得其積而後求其長短廣狹斜正之形，有非乘除所能盡者，故必以商除之。然而商除亦不能盡也，而又立正負廉隅之法以增損附益之，故其為術也難。

\switchcolumn

In Li Ye's original work, he first presents the \textbf{gougu zong tu} (general diagram of base, altitude, and hypotenuse), using characters such as Heaven, Earth, Sun, Moon, Mountains, Rivers, Dawn, Dusk, Vermilion, Azure, Overflowing, and Spring, as well as the names of the Four Directions, Five Phases, and Eight Trigrams to mark the intersections in the diagram in detail. The various shapes formed at these intersections serve to communicate the principles of commerce and works. To seek area from shape pertains to methods like rectangular fields and works; to seek shape from area pertains to methods like lesser breadth and right triangles. Seeking area from shape---first obtaining the shape and then finding its area---is therefore an easy method. But seeking shape from area---first obtaining the area and then finding its dimensions of length, width, and slant---involves cases that multiplication and division alone cannot fully resolve, so one must use root extraction. Yet even root extraction cannot exhaust all cases; thus the methods of positive and negative edges and corners were established to adjust and supplement them, making this a difficult art.

\switchcolumn

余自幼好習數學，晚得荊川唐太史所錄\textbf{測圓海鏡}一書，乃元翰林學士欒城李公冶所著。雖專主於勾股求容圓容方一術，然其中間如平方、立方、三乘方、帶從、減從、益廉、減廉、正隅、負隅諸法，凡所謂以積求形者，皆盡之矣。但其每條下細草，俱徑立天元一，反覆合之，而無下手之術。使後學之士，茫然無門路可入，輒不自揆，每章去其細草，立一算術。又以其所立通勾邊股之屬各以類分之，語義稍繁者略加芟損，名曰\textbf{測圓海鏡分類釋術}。匪敢僭改前賢著述，惟以便下學云爾。

\switchcolumn

I have loved the study of mathematics since youth. In my later years I obtained a copy of the \textit{Ceyuan Haijing} recorded by the Grand Historian Tang of Jingchuan---a work composed by the Hanlin Academician Lord Li Ye of Luancheng during the Yuan Dynasty. Although it specializes chiefly in the method of finding the inscribed and escribed circles of right triangles, it also covers square root, cube root, triple multiplication root extraction, carrying-associate, subtracting-associate, adding-edge, subtracting-edge, positive-corner, and negative-corner methods---in short, all the so-called ``seeking shape from area'' methods are fully treated. However, beneath each entry, the detailed working (\textit{xi cao}) all proceeded directly by establishing the celestial element unity (\textit{tian yuan yi}), folding and unfolding it repeatedly, without showing the method by which to begin. This left later students lost, with no door or path by which to enter. Unable to assess my own capabilities, I removed the \textit{xi cao} from each chapter and established an explicit arithmetic procedure (\textit{suan shu}). I also classified the various quantities such as \textit{tong gou} (universal base) and \textit{bian gu} (edge altitude) by category, slightly trimming passages that were overly verbose, and named the result \textit{Ceyuan Haijing Fenlei Shishu}. I dare not claim to have revised the works of past sages; my sole aim is to facilitate beginners.

\switchcolumn

今夫世之論數者，俱視為末藝，故高明者不屑為之，而執泥者遂以為占驗之法。雖欒城公自序亦以為九九賤伎，殊不知君子之學，自性命道德之外皆藝也。與其徒費精神於佔畢之間，又不若留情於此。不惟可以取樂，亦足以為養心之助焉。後之有同此好者，當以余言為然否耶？

\switchcolumn

Now, those who discuss numbers in the world today all regard it as a minor art; therefore the brilliant disdain to practice it, while the stubborn take it as a method of divination. Even Lord Li Ye of Luancheng himself, in his preface, considered the multiplication table a lowly craft. Little do they realize that beyond the study of human nature and moral principles, all is art. Rather than wasting one's spirit on rote recitation, it would be better to devote attention to this subject. Not only can it provide pleasure, but it also serves to cultivate the mind. If any later generations share this fondness, will they find my words to be true or not?

\switchcolumn

\begin{flushright}
嘉靖庚戌夏五月朔\\
吳興顧應祥誌
\end{flushright}

\switchcolumn

\begin{flushright}
First day of the fifth month, summer,\\
of the Gengxu year of Jiajing (1550)\\
Inscribed by Gu Yingxiang of Wuxing
\end{flushright}

\end{paracol}

\clearpage

%% ========================================================================
%% PREFACE (原序)
%% ========================================================================
\section*{原序 / Original Preface}
\addcontentsline{toc}{section}{原序 / Original Preface}
\markboth{原序 / Original Preface}{原序 / Original Preface}

\begin{paracol}{2}

\noindent\textbf{原序}

\medskip

天地之所以神變化而生萬物者，陰陽而已。一陰一陽，交互錯綜，而變化無窮焉。聖人因其交互錯綜之不齊，而置為數術以測之。於是乎天地之高深，日月之出沒，鬼神之幽秘，皆可得而知之矣。然數之為術，雖千變萬化之不同，而其要不過一開一闔而已。開者除也，闔者乘也。而又有以形求積、以積求形之異。古之為數者，有九九者，其用也。是故用之以貿易則為粟米，用之以分別差等較量遠近，則為差分、為均輸。因其末而欲知其本，求其故，千歲之日至，可坐而致也。跡其所以求天與星辰之高遠，非勾股何以御之？而周官土圭測景之術，所以窮地之塊扎無垠，若指諸掌，亦此術也。豈得以為古奧而棄之不講乎？

\switchcolumn

\noindent\textbf{Original Preface}

\medskip

That by which Heaven and Earth effect divine transformations and give birth to the myriad things is none other than Yin and Yang. One Yin and one Yang, interweaving and interlocking, produce inexhaustible transformations. The Sages, because of the irregularities of this interweaving and interlocking, established mathematical arts to measure them. Thus the heights and depths of Heaven and Earth, the risings and settings of the Sun and Moon, and the mysteries of spirits and ghosts can all be known. Yet the art of numbers, though it has a thousand transformations and ten thousand variations, is at root nothing more than opening and closing. Opening is division; closing is multiplication. Furthermore, there is the distinction between seeking area from shape and seeking shape from area. The ancients who practiced mathematics had the multiplication table as their tool. Applied to trade it becomes grain measures; applied to distinguishing gradations and comparing distances it becomes difference distribution and equal transport. Following the branch to know the root, seeking the cause---the winter solstice of a thousand years hence can be found while seated. Tracing the method by which the heights of Heaven and the stars are sought, what besides the right triangle could manage it? And the method of the Zhou officials who measured shadows with the gnomon, thereby fathoming the boundless extent of the Earth as if pointing to their palms---this too is the same art. How can one regard it as archaic and abandon its study?

\switchcolumn

此公語愚分釋此書之盛心也，因請於公錄一帙而刻之，梓播之四方，傳之後世，以見公之經濟不獨惠我滇雲，而推明朕兆根極領要，以繼絕學。俟後賢者尚有考於斯焉。

\switchcolumn

This was the generous intention of the gentleman in speaking to me about classifying and explaining this book. I therefore requested from him a copy to engrave and print, to disseminate in all directions and transmit to later generations, so that the gentleman's practical wisdom may benefit not only our land of Dianyun, but also illuminate the fundamental principles and essential points, thereby continuing a lost learning. I await the wise men of the future who may still examine this work.

\switchcolumn

\begin{flushright}
嘉靖庚戌夏五月朔日\\
古濠沐朝弼謹序
\end{flushright}

\switchcolumn

\begin{flushright}
First day of the fifth month, summer,\\
of the Gengxu year of Jiajing (1550)\\
Respectfully preface by Mu Chaobi of Guhao
\end{flushright}

\end{paracol}

\clearpage

%% ========================================================================
%% VOLUME 1 (卷一)
%% ========================================================================
\section{卷一 / Volume 1}
\markboth{卷一 / Volume 1}{卷一 / Volume 1}

\begin{paracol}{2}

\noindent\textbf{卷一}

\smallskip

卷一專論\textbf{通勾股求容圓}之法，凡十問。通勾股者，以通勾、通股、通弦三者之關係，求圓城之徑也。

\switchcolumn

\noindent\textbf{Volume 1}

\smallskip

Volume 1 is devoted to the method of \textbf{finding the inscribed circle using the universal right triangle} (tong gougu qiu rong yuan), comprising ten problems. The universal right triangle uses the relationships among the universal base (tong gou), universal altitude (tong gu), and universal hypotenuse (tong xian) to find the diameter of the circular city.

\end{paracol}

\bigskip

%% ========================================================================
%% Chapter 1: Tong gougu rong yuan
%% ========================================================================
\subsection*{通勾股求容圓 / Finding the Inscribed Circle via the Universal Right Triangle}
\addcontentsline{toc}{subsection}{通勾股求容圓 / Universal Right Triangle}

\begin{paracol}{2}

\noindent\textbf{總論：} 通勾股求容圓者，以勾股相乘倍之為實，勾股和為法，除之得圓徑。其術之要，在於以通勾、通股之關係，推求容圓之徑也。

\switchcolumn

\noindent\textbf{General Discussion:} To find the inscribed circle via the universal right triangle, multiply the base by the altitude, double it to make the dividend; take the sum of base and altitude as the divisor; dividing gives the circle's diameter. The essence of the method lies in using the relationship between the universal base and universal altitude to derive the diameter of the inscribed circle.

\end{paracol}

\bigskip

\begin{paracol}{2}

\noindent\textbf{第1題 通勾股求容圓一}

\wen{甲南行六百步，望乙與城相參直。問城徑。}
\da{城徑二百四十步。}
\shu{勾股相乘倍之為實，勾股求弦併勾股為弦和。弦和較即容圓徑也。}
\shi{此勾股求容圓徑也。東行為通勾，南行為通股。以通勾股求通弦和，較弦和較即容圓徑也。}

\switchcolumn

\noindent\textbf{Problem 1: Inscribed Circle via Universal Right Triangle, 1}

\wenE{A walks south 600 paces and sees B aligned with the city [edge]. Find the city's diameter.}
\daE{The city diameter is 240 paces.}
\shuE{Multiply the base by the altitude and double it to make the dividend; find the hypotenuse from base and altitude, combine base and altitude with the hypotenuse to make the hypotenuse-sum; the difference between hypotenuse and sum is the inscribed circle's diameter.}
\shiE{This is the method of finding the inscribed circle diameter via the right triangle. The eastward walk is the universal base; the southward walk is the universal altitude. Using the universal base and altitude to find the universal hypotenuse-sum, the difference between the hypotenuse and sum gives the inscribed circle's diameter.}

\end{paracol}

\bigskip

\begin{paracol}{2}

\noindent\textbf{第2題 通勾股求容圓二}

\wen{甲乙二人俱在城西門。甲南行四百八十步，乙穿城東行二百五十步見之。問城徑。}
\da{城徑二百四十步。}
\shu{勾股相乘倍之為實，勾股求弦併勾股為弦和。弦和較即容圓徑也。}
\shi{此勾上容圓也。南行為邊股，東行為邊勾也。以邊勾股求容圓，與通勾股求通圓同法。}

\switchcolumn

\noindent\textbf{Problem 2: Inscribed Circle via Universal Right Triangle, 2}

\wenE{A and B are both at the city's west gate. A walks south 480 paces; B goes through the city walking east 250 paces and sees A. Find the city's diameter.}
\daE{The city diameter is 240 paces.}
\shuE{Multiply the base by the altitude and double it to make the dividend; find the hypotenuse from base and altitude, combine base and altitude with the hypotenuse to make the hypotenuse-sum; the difference between hypotenuse and sum is the inscribed circle's diameter.}
\shiE{This is finding the inscribed circle on the base. The southward walk is the edge-altitude; the eastward walk is the edge-base. Using the edge-base and edge-altitude to find the inscribed circle follows the same method as using the universal base and altitude to find the universal circle.}

\end{paracol}

\bigskip

\begin{paracol}{2}

\noindent\textbf{第3題 通勾股求容圓三}

\wen{甲出北門東行，乙出西門南行，相望見之。甲東行二百步，乙南行四百二十五步。問城徑。}
\da{城徑二百四十步。}
\shu{勾股相乘倍之為實，勾股求弦併勾股為弦和。弦和較即容圓徑也。}
\shi{此亦邊勾股求容圓也。以邊勾股之數，約之得圓徑。}

\switchcolumn

\noindent\textbf{Problem 3: Inscribed Circle via Universal Right Triangle, 3}

\wenE{A exits the north gate and walks east; B exits the west gate and walks south. They see each other. A walks east 200 paces; B walks south 425 paces. Find the city's diameter.}
\daE{The city diameter is 240 paces.}
\shuE{Multiply the base by the altitude and double it to make the dividend; find the hypotenuse from base and altitude, combine base and altitude with the hypotenuse to make the hypotenuse-sum; the difference between hypotenuse and sum is the inscribed circle's diameter.}
\shiE{This too is finding the inscribed circle via the edge right triangle. Using the edge-base and edge-altitude values and reducing them gives the circle's diameter.}

\end{paracol}

\bigskip

\begin{paracol}{2}

\noindent\textbf{第4題 通勾股求容圓四}

\wen{甲乙二人俱在城中心立。乙穿城東行一百三十六步，甲出南門東行二百步見之。問城徑。}
\da{城徑二百四十步。}
\shu{勾股相乘倍之為實，勾股求弦併勾股為弦和。弦和較即容圓徑也。}
\shi{此皇極勾股求容圓也。以通勾股之半為皇極勾股，其容圓徑之半即城半徑也。}

\switchcolumn

\noindent\textbf{Problem 4: Inscribed Circle via Universal Right Triangle, 4}

\wenE{A and B both stand at the center of the city. B goes through the city walking east 136 paces; A exits the south gate and walks east 200 paces and sees B. Find the city's diameter.}
\daE{The city diameter is 240 paces.}
\shuE{Multiply the base by the altitude and double it to make the dividend; find the hypotenuse from base and altitude, combine base and altitude with the hypotenuse to make the hypotenuse-sum; the difference between hypotenuse and sum is the inscribed circle's diameter.}
\shiE{This is finding the inscribed circle via the supreme-pole right triangle. Taking half the universal base and altitude makes the supreme-pole base and altitude; half of its inscribed circle diameter is the city's radius.}

\end{paracol}

\bigskip

\begin{paracol}{2}

\noindent\textbf{第5題 通勾股求容圓五}

\wen{甲出南門東行，乙出東門南行，相望見之。甲東行七十二步，乙南行三十步。問城徑。}
\da{城徑二百四十步。}
\shu{勾股相乘倍之為實，勾股求弦併勾股為弦和。弦和較即容圓徑也。}
\shi{此太虛勾股求容圓也。以勾股之數約之得容圓徑。}

\switchcolumn

\noindent\textbf{Problem 5: Inscribed Circle via Universal Right Triangle, 5}

\wenE{A exits the south gate and walks east; B exits the east gate and walks south. They see each other. A walks east 72 paces; B walks south 30 paces. Find the city's diameter.}
\daE{The city diameter is 240 paces.}
\shuE{Multiply the base by the altitude and double it to make the dividend; find the hypotenuse from base and altitude, combine base and altitude with the hypotenuse to make the hypotenuse-sum; the difference between hypotenuse and sum is the inscribed circle's diameter.}
\shiE{This is finding the inscribed circle via the grand-void right triangle. Using the base and altitude values and reducing them gives the inscribed circle diameter.}

\end{paracol}

\bigskip

\begin{paracol}{2}

\noindent\textbf{第6題 通勾股求容圓六}

\wen{甲出西門南行，乙出南門東行，相望見之。甲南行八十步，乙東行五十一步。問城徑。}
\da{城徑二百四十步。}
\shu{勾股相乘倍之為實，勾股求弦併勾股為弦和。弦和較即容圓徑也。}
\shi{此明勾股求容圓也。以明勾股之數，約之得容圓徑。}

\switchcolumn

\noindent\textbf{Problem 6: Inscribed Circle via Universal Right Triangle, 6}

\wenE{A exits the west gate and walks south; B exits the south gate and walks east. They see each other. A walks south 80 paces; B walks east 51 paces. Find the city's diameter.}
\daE{The city diameter is 240 paces.}
\shuE{Multiply the base by the altitude and double it to make the dividend; find the hypotenuse from base and altitude, combine base and altitude with the hypotenuse to make the hypotenuse-sum; the difference between hypotenuse and sum is the inscribed circle's diameter.}
\shiE{This is finding the inscribed circle via the bright right triangle. Using the bright base and altitude values and reducing them gives the inscribed circle diameter.}

\end{paracol}

\bigskip

\begin{paracol}{2}

\noindent\textbf{第7題 通勾股求容圓七}

\wen{甲出東門北行，乙出北門西行，相望見之。甲北行三十步，乙西行一十六步。問城徑。}
\da{城徑二百四十步。}
\shu{勾股相乘倍之為實，勾股求弦併勾股為弦和。弦和較即容圓徑也。}
\shi{此重勾股求容圓也。以重勾股之數，約之得容圓徑。}

\switchcolumn

\noindent\textbf{Problem 7: Inscribed Circle via Universal Right Triangle, 7}

\wenE{A exits the east gate and walks north; B exits the north gate and walks west. They see each other. A walks north 30 paces; B walks west 16 paces. Find the city's diameter.}
\daE{The city diameter is 240 paces.}
\shuE{Multiply the base by the altitude and double it to make the dividend; find the hypotenuse from base and altitude, combine base and altitude with the hypotenuse to make the hypotenuse-sum; the difference between hypotenuse and sum is the inscribed circle's diameter.}
\shiE{This is finding the inscribed circle via the double right triangle. Using the double base and altitude values and reducing them gives the inscribed circle diameter.}

\end{paracol}

\bigskip

\begin{paracol}{2}

\noindent\textbf{第8題 通勾股求容圓八}

\wen{甲出東門直行，乙出南門直行，相望見之。甲東行三十步，乙南行十六步。問城徑。}
\da{城徑二百四十步。}
\shu{勾股相乘倍之為實，勾股求弦併勾股為弦和。弦和較即容圓徑也。}
\shi{此虛勾股求容圓也。以虛勾股之數約之，得圓徑。}

\switchcolumn

\noindent\textbf{Problem 8: Inscribed Circle via Universal Right Triangle, 8}

\wenE{A exits the east gate and walks straight; B exits the south gate and walks straight. They see each other. A walks east 30 paces; B walks south 16 paces. Find the city's diameter.}
\daE{The city diameter is 240 paces.}
\shuE{Multiply the base by the altitude and double it to make the dividend; find the hypotenuse from base and altitude, combine base and altitude with the hypotenuse to make the hypotenuse-sum; the difference between hypotenuse and sum is the inscribed circle's diameter.}
\shiE{This is finding the inscribed circle via the void right triangle. Using the void base and altitude values and reducing them gives the circle's diameter.}

\end{paracol}

\bigskip

\begin{paracol}{2}

\noindent\textbf{第9題 通勾股求容圓九}

\wen{甲出南門直行，乙出東門直行，相望見之。甲南行四十八步，乙東行二十步。問城徑。}
\da{城徑二百四十步。}
\shu{勾股相乘倍之為實，勾股求弦併勾股為弦和。弦和較即容圓徑也。}
\shi{此以勾股之數求容圓也。以勾股之比例推之，得圓徑。}

\switchcolumn

\noindent\textbf{Problem 9: Inscribed Circle via Universal Right Triangle, 9}

\wenE{A exits the south gate and walks straight; B exits the east gate and walks straight. They see each other. A walks south 48 paces; B walks east 20 paces. Find the city's diameter.}
\daE{The city diameter is 240 paces.}
\shuE{Multiply the base by the altitude and double it to make the dividend; find the hypotenuse from base and altitude, combine base and altitude with the hypotenuse to make the hypotenuse-sum; the difference between hypotenuse and sum is the inscribed circle's diameter.}
\shiE{This uses the base and altitude values to find the inscribed circle. Using the ratio of base to altitude and extending it gives the circle's diameter.}

\end{paracol}

\bigskip

\begin{paracol}{2}

\noindent\textbf{第10題 通勾股求容圓十}

\wen{甲出西門直行，乙出北門直行，相望見之。甲西行二百步，乙北行六十步。問城徑。}
\da{城徑二百四十步。}
\shu{勾股相乘倍之為實，勾股求弦併勾股為弦和。弦和較即容圓徑也。}
\shi{此以大差勾股求容圓也。以大差勾股之數約之，得圓徑。}

\switchcolumn

\noindent\textbf{Problem 10: Inscribed Circle via Universal Right Triangle, 10}

\wenE{A exits the west gate and walks straight; B exits the north gate and walks straight. They see each other. A walks west 200 paces; B walks north 60 paces. Find the city's diameter.}
\daE{The city diameter is 240 paces.}
\shuE{Multiply the base by the altitude and double it to make the dividend; find the hypotenuse from base and altitude, combine base and altitude with the hypotenuse to make the hypotenuse-sum; the difference between hypotenuse and sum is the inscribed circle's diameter.}
\shiE{This is finding the inscribed circle via the great-difference right triangle. Using the great-difference base and altitude values and reducing them gives the circle's diameter.}

\end{paracol}

\clearpage

%% ========================================================================
%% Chapter 2: Bian gougu rong yuan
%% ========================================================================
\subsection*{邊勾股求容圓 / Finding the Inscribed Circle via the Edge Right Triangle}
\addcontentsline{toc}{subsection}{邊勾股求容圓 / Edge Right Triangle}

\begin{paracol}{2}

\noindent\textbf{總論：} 邊勾股求容圓者，以邊勾、邊股之數，用通勾股之法，求容圓之徑也。

\switchcolumn

\noindent\textbf{General Discussion:} To find the inscribed circle via the edge right triangle, use the edge-base and edge-altitude values, applying the method of the universal right triangle, to find the diameter of the inscribed circle.

\end{paracol}

\bigskip

\begin{paracol}{2}

\noindent\textbf{第11題 邊勾股求容圓一}

\wen{甲出北門東行二百步，乙出東門南行一百二十步，相望見之。問城徑。}
\da{城徑二百四十步。}
\shu{邊勾股相乘倍之為實，邊勾股求弦併邊勾股為弦和。弦和較即容圓徑也。}
\shi{此邊勾股求容圓也。邊勾股者，通勾股之半也。以邊勾股之數求容圓，與通勾股同法。}

\switchcolumn

\noindent\textbf{Problem 11: Inscribed Circle via Edge Right Triangle, 1}

\wenE{A exits the north gate and walks east 200 paces; B exits the east gate and walks south 120 paces. They see each other. Find the city's diameter.}
\daE{The city diameter is 240 paces.}
\shuE{Multiply the edge-base by the edge-altitude and double it to make the dividend; find the hypotenuse from the edge-base and edge-altitude, combine them to make the hypotenuse-sum; the difference between hypotenuse and sum is the inscribed circle's diameter.}
\shiE{This is finding the inscribed circle via the edge right triangle. The edge-base and edge-altitude are half of the universal base and altitude. Using the edge-base and edge-altitude to find the inscribed circle follows the same method as the universal right triangle.}

\end{paracol}

\bigskip

\begin{paracol}{2}

\noindent\textbf{第12題 邊勾股求容圓二}

\wen{甲乙二人俱在城西門。甲南行四百八十步，乙穿城東行二百五十步見之。問城徑。}
\da{城徑二百四十步。}
\shu{勾股相乘倍之為實，勾股求弦併勾股為弦和。弦和較即容圓徑也。}
\shi{此勾上容圓也。南行為邊股，東行為邊勾也。以邊勾股求容圓，與通勾股求通圓同法。}

\switchcolumn

\noindent\textbf{Problem 12: Inscribed Circle via Edge Right Triangle, 2}

\wenE{A and B are both at the city's west gate. A walks south 480 paces; B goes through the city walking east 250 paces and sees A. Find the city's diameter.}
\daE{The city diameter is 240 paces.}
\shuE{Multiply the base by the altitude and double it to make the dividend; find the hypotenuse from base and altitude, combine base and altitude with the hypotenuse to make the hypotenuse-sum; the difference between hypotenuse and sum is the inscribed circle's diameter.}
\shiE{This is finding the inscribed circle on the base. The southward walk is the edge-altitude; the eastward walk is the edge-base. Using the edge-base and edge-altitude to find the inscribed circle follows the same method as the universal right triangle.}

\end{paracol}

\bigskip

\begin{paracol}{2}

\noindent\textbf{第13題 邊勾股求容圓三}

\wen{甲出北門東行二百步，乙出東門南行一百二十步，相望見之。問城徑。}
\da{城徑二百四十步。}
\shu{邊勾股相乘倍之為實，邊勾股求弦併邊勾股為弦和。弦和較即容圓徑也。}
\shi{此邊勾股求容圓也。邊勾股者，通勾股之半也。以邊勾股之數求容圓，與通勾股同法。}

\switchcolumn

\noindent\textbf{Problem 13: Inscribed Circle via Edge Right Triangle, 3}

\wenE{A exits the north gate and walks east 200 paces; B exits the east gate and walks south 120 paces. They see each other. Find the city's diameter.}
\daE{The city diameter is 240 paces.}
\shuE{Multiply the edge-base by the edge-altitude and double it to make the dividend; find the hypotenuse from the edge-base and edge-altitude, combine them to make the hypotenuse-sum; the difference between hypotenuse and sum is the inscribed circle's diameter.}
\shiE{This is finding the inscribed circle via the edge right triangle. The edge-base and edge-altitude are half of the universal base and altitude. Using the edge-base and edge-altitude to find the inscribed circle follows the same method as the universal right triangle.}

\end{paracol}

\bigskip

\begin{paracol}{2}

\noindent\textbf{第14題 邊勾股求容圓四}

\wen{甲出南門東行一百二十步，乙出東門南行五十步，相望見之。問城徑。}
\da{城徑二百四十步。}
\shu{邊勾股相乘倍之為實，邊勾股求弦併邊勾股為弦和。弦和較即容圓徑也。}
\shi{此邊勾股求容圓也。以邊勾股之數求容圓，與通勾股同法。}

\switchcolumn

\noindent\textbf{Problem 14: Inscribed Circle via Edge Right Triangle, 4}

\wenE{A exits the south gate and walks east 120 paces; B exits the east gate and walks south 50 paces. They see each other. Find the city's diameter.}
\daE{The city diameter is 240 paces.}
\shuE{Multiply the edge-base by the edge-altitude and double it to make the dividend; find the hypotenuse from the edge-base and edge-altitude, combine them to make the hypotenuse-sum; the difference between hypotenuse and sum is the inscribed circle's diameter.}
\shiE{This is finding the inscribed circle via the edge right triangle. Using the edge-base and edge-altitude values to find the inscribed circle follows the same method as the universal right triangle.}

\end{paracol}

\bigskip

\begin{paracol}{2}

\noindent\textbf{第15題 邊勾股求容圓五}

\wen{甲出西門南行一百步，乙出南門西行八十步，相望見之。問城徑。}
\da{城徑二百四十步。}
\shu{邊勾股相乘倍之為實，邊勾股求弦併邊勾股為弦和。弦和較即容圓徑也。}
\shi{此邊勾股求容圓也。以邊勾股之數求容圓，與通勾股同法。}

\switchcolumn

\noindent\textbf{Problem 15: Inscribed Circle via Edge Right Triangle, 5}

\wenE{A exits the west gate and walks south 100 paces; B exits the south gate and walks west 80 paces. They see each other. Find the city's diameter.}
\daE{The city diameter is 240 paces.}
\shuE{Multiply the edge-base by the edge-altitude and double it to make the dividend; find the hypotenuse from the edge-base and edge-altitude, combine them to make the hypotenuse-sum; the difference between hypotenuse and sum is the inscribed circle's diameter.}
\shiE{This is finding the inscribed circle via the edge right triangle. Using the edge-base and edge-altitude values to find the inscribed circle follows the same method as the universal right triangle.}

\end{paracol}

\clearpage

%% ========================================================================
%% Chapter 3: Di gu xian qiu tong yuanjing
%% ========================================================================
\subsection*{底股弦求通圓徑 / Finding the Universal Circle Diameter via the Base Altitude and Hypotenuse}
\addcontentsline{toc}{subsection}{底股弦求通圓徑 / Base Altitude and Hypotenuse}

\begin{paracol}{2}

\noindent\textbf{總論：} 底股弦求通圓徑者，以底股、底弦之數，用弦較之法，求通圓之徑也。

\switchcolumn

\noindent\textbf{General Discussion:} To find the universal circle diameter via the base altitude and hypotenuse, use the base-altitude and base-hypotenuse values, applying the hypotenuse-difference method, to find the diameter of the universal circle.

\end{paracol}

\bigskip

\begin{paracol}{2}

\noindent\textbf{第16題 底股弦求通圓徑一}

\wen{問底股弦求通圓徑。}
\da{通圓徑二百四十步。}
\shu{弦昇減股昇開其餘，得勾如前法求之。}
\shi{此以底股弦之關係，求通圓徑也。底股弦者，大勾股弦之屬也。}

\switchcolumn

\noindent\textbf{Problem 16: Universal Circle Diameter via Base Altitude and Hypotenuse, 1}

\wenE{Question: Find the universal circle diameter using the base altitude and hypotenuse.}
\daE{The universal circle diameter is 240 paces.}
\shuE{Raise the hypotenuse to the power, subtract the altitude raised to the power, and open the remainder to obtain the base; then proceed as in the previous method.}
\shiE{This uses the relationship between the base-altitude and base-hypotenuse to find the universal circle diameter. The base-altitude and base-hypotenuse belong to the greater right triangle category.}

\end{paracol}

\bigskip

\begin{paracol}{2}

\noindent\textbf{第17題 底股弦求通圓徑二}

\wen{甲出北門東行，乙出東門南行，相望見之。問以底股弦求通圓徑。}
\da{通圓徑二百四十步。}
\shu{弦昇減股昇開其餘，得勾如前法求之。}
\shi{此底股弦求通圓徑也。以弦昇減股昇開其餘得勾，再以勾股求容圓。}

\switchcolumn

\noindent\textbf{Problem 17: Universal Circle Diameter via Base Altitude and Hypotenuse, 2}

\wenE{A exits the north gate and walks east; B exits the east gate and walks south. They see each other. Find the universal circle diameter using the base altitude and hypotenuse.}
\daE{The universal circle diameter is 240 paces.}
\shuE{Raise the hypotenuse to the power, subtract the altitude raised to the power, and open the remainder to obtain the base; then proceed as in the previous method.}
\shiE{This is finding the universal circle diameter via the base altitude and hypotenuse. Subtracting the altitude-power from the hypotenuse-power and opening the remainder gives the base; then use base and altitude to find the inscribed circle.}

\end{paracol}

\bigskip

\begin{paracol}{2}

\noindent\textbf{第18題 底股弦求通圓徑三}

\wen{甲出南門直行，乙出東門直行，相望見之。問以底股弦求通圓徑。}
\da{通圓徑二百四十步。}
\shu{弦昇減股昇開其餘，得勾如前法求之。}
\shi{此底股弦求通圓徑也。以弦昇減股昇開其餘得勾，再以勾股求容圓。}

\switchcolumn

\noindent\textbf{Problem 18: Universal Circle Diameter via Base Altitude and Hypotenuse, 3}

\wenE{A exits the south gate and walks straight; B exits the east gate and walks straight. They see each other. Find the universal circle diameter using the base altitude and hypotenuse.}
\daE{The universal circle diameter is 240 paces.}
\shuE{Raise the hypotenuse to the power, subtract the altitude raised to the power, and open the remainder to obtain the base; then proceed as in the previous method.}
\shiE{This is finding the universal circle diameter via the base altitude and hypotenuse. Subtracting the altitude-power from the hypotenuse-power and opening the remainder gives the base; then use base and altitude to find the inscribed circle.}

\end{paracol}

\clearpage

%% ========================================================================
%% Chapter 4: Huangji gougu qiu rong yuan
%% ========================================================================
\subsection*{皇極勾股求容圓 / Finding the Inscribed Circle via the Supreme-Pole Right Triangle}
\addcontentsline{toc}{subsection}{皇極勾股求容圓 / Supreme-Pole Right Triangle}

\begin{paracol}{2}

\noindent\textbf{總論：} 皇極勾股求容圓者，以皇極勾、皇極股之數，用勾股容圓之法，求圓徑也。皇極勾股者，通勾股之半也。

\switchcolumn

\noindent\textbf{General Discussion:} To find the inscribed circle via the supreme-pole right triangle, use the supreme-pole base and supreme-pole altitude values, applying the right-triangle inscribed-circle method, to find the circle's diameter. The supreme-pole right triangle is half of the universal right triangle.

\end{paracol}

\bigskip

\begin{paracol}{2}

\noindent\textbf{第19題 皇極勾股求容圓一}

\wen{甲乙二人俱在城中心立。乙穿城東行一百三十六步，甲出南門東行二百步見之。問城徑。}
\da{城徑二百四十步。}
\shu{皇極勾股相乘倍之為實，皇極勾股和為法，除之得皇極容圓徑，倍之為城徑。}
\shi{此皇極勾股求容圓也。以皇極勾股之半為通勾股之半，其容圓徑之半即城半徑也。}

\switchcolumn

\noindent\textbf{Problem 19: Inscribed Circle via Supreme-Pole Right Triangle, 1}

\wenE{A and B both stand at the center of the city. B goes through the city walking east 136 paces; A exits the south gate and walks east 200 paces and sees B. Find the city's diameter.}
\daE{The city diameter is 240 paces.}
\shuE{Multiply the supreme-pole base by the supreme-pole altitude and double it to make the dividend; the sum of supreme-pole base and altitude is the divisor; dividing gives the supreme-pole inscribed circle diameter; doubling it gives the city diameter.}
\shiE{This is finding the inscribed circle via the supreme-pole right triangle. Taking half the supreme-pole base and altitude gives half the universal base and altitude; half of the inscribed circle diameter is the city's radius.}

\end{paracol}

\bigskip

\begin{paracol}{2}

\noindent\textbf{第20題 皇極勾股求容圓二}

\wen{甲出北門東行一百步，乙出東門南行八十步，相望見之。問城徑。}
\da{城徑二百四十步。}
\shu{皇極勾股相乘倍之為實，皇極勾股和為法，除之得皇極容圓徑，倍之為城徑。}
\shi{此皇極勾股求容圓也。以皇極勾股之數，約之得圓徑。}

\switchcolumn

\noindent\textbf{Problem 20: Inscribed Circle via Supreme-Pole Right Triangle, 2}

\wenE{A exits the north gate and walks east 100 paces; B exits the east gate and walks south 80 paces. They see each other. Find the city's diameter.}
\daE{The city diameter is 240 paces.}
\shuE{Multiply the supreme-pole base by the supreme-pole altitude and double it to make the dividend; the sum of supreme-pole base and altitude is the divisor; dividing gives the supreme-pole inscribed circle diameter; doubling it gives the city diameter.}
\shiE{This is finding the inscribed circle via the supreme-pole right triangle. Using the supreme-pole base and altitude values and reducing them gives the circle's diameter.}

\end{paracol}

\bigskip

\begin{paracol}{2}

\noindent\textbf{第21題 皇極勾股求容圓三}

\wen{甲出南門東行一百五十步，乙出東門南行一百步，相望見之。問城徑。}
\da{城徑二百四十步。}
\shu{皇極勾股相乘倍之為實，皇極勾股和為法，除之得皇極容圓徑，倍之為城徑。}
\shi{此皇極勾股求容圓也。以皇極勾股之數，約之得圓徑。}

\switchcolumn

\noindent\textbf{Problem 21: Inscribed Circle via Supreme-Pole Right Triangle, 3}

\wenE{A exits the south gate and walks east 150 paces; B exits the east gate and walks south 100 paces. They see each other. Find the city's diameter.}
\daE{The city diameter is 240 paces.}
\shuE{Multiply the supreme-pole base by the supreme-pole altitude and double it to make the dividend; the sum of supreme-pole base and altitude is the divisor; dividing gives the supreme-pole inscribed circle diameter; doubling it gives the city diameter.}
\shiE{This is finding the inscribed circle via the supreme-pole right triangle. Using the supreme-pole base and altitude values and reducing them gives the circle's diameter.}

\end{paracol}

\bigskip

\begin{paracol}{2}

\noindent\textbf{第22題 皇極勾股求容圓四}

\wen{甲出西門南行二百步，乙出南門西行一百五十步，相望見之。問城徑。}
\da{城徑二百四十步。}
\shu{皇極勾股相乘倍之為實，皇極勾股和為法，除之得皇極容圓徑，倍之為城徑。}
\shi{此皇極勾股求容圓也。以皇極勾股之數，約之得圓徑。}

\switchcolumn

\noindent\textbf{Problem 22: Inscribed Circle via Supreme-Pole Right Triangle, 4}

\wenE{A exits the west gate and walks south 200 paces; B exits the south gate and walks west 150 paces. They see each other. Find the city's diameter.}
\daE{The city diameter is 240 paces.}
\shuE{Multiply the supreme-pole base by the supreme-pole altitude and double it to make the dividend; the sum of supreme-pole base and altitude is the divisor; dividing gives the supreme-pole inscribed circle diameter; doubling it gives the city diameter.}
\shiE{This is finding the inscribed circle via the supreme-pole right triangle. Using the supreme-pole base and altitude values and reducing them gives the circle's diameter.}

\end{paracol}

\clearpage

%% ========================================================================
%% Chapter 5: Taiyin gougu qiu rong yuan
%% ========================================================================
\subsection*{太陰勾股求容圓 / Finding the Inscribed Circle via the Grand-Yin Right Triangle}
\addcontentsline{toc}{subsection}{太陰勾股求容圓 / Grand-Yin Right Triangle}

\begin{paracol}{2}

\noindent\textbf{總論：} 太陰勾股求容圓者，以太陰勾、太陰股之數，用勾股容圓之法，求圓徑也。

\switchcolumn

\noindent\textbf{General Discussion:} To find the inscribed circle via the grand-yin right triangle, use the grand-yin base and grand-yin altitude values, applying the right-triangle inscribed-circle method, to find the circle's diameter.

\end{paracol}

\bigskip

\begin{paracol}{2}

\noindent\textbf{第23題 太陰勾股求容圓一}

\wen{甲出南門東行七十二步，乙出東門南行三十步，相望見之。問城徑。}
\da{城徑二百四十步。}
\shu{太陰勾股相乘倍之為實，太陰勾股和為法，除之得太陰容圓徑，以比例推之得城徑。}
\shi{此太陰勾股求容圓也。以太陰勾股之數，約之得圓徑。}

\switchcolumn

\noindent\textbf{Problem 23: Inscribed Circle via Grand-Yin Right Triangle, 1}

\wenE{A exits the south gate and walks east 72 paces; B exits the east gate and walks south 30 paces. They see each other. Find the city's diameter.}
\daE{The city diameter is 240 paces.}
\shuE{Multiply the grand-yin base by the grand-yin altitude and double it to make the dividend; the sum of grand-yin base and altitude is the divisor; dividing gives the grand-yin inscribed circle diameter; extending by proportion gives the city diameter.}
\shiE{This is finding the inscribed circle via the grand-yin right triangle. Using the grand-yin base and altitude values and reducing them gives the circle's diameter.}

\end{paracol}

\bigskip

\begin{paracol}{2}

\noindent\textbf{第24題 太陰勾股求容圓二}

\wen{甲出北門東行五十步，乙出東門北行二十步，相望見之。問城徑。}
\da{城徑二百四十步。}
\shu{太陰勾股相乘倍之為實，太陰勾股和為法，除之得太陰容圓徑，以比例推之得城徑。}
\shi{此太陰勾股求容圓也。以太陰勾股之數，約之得圓徑。}

\switchcolumn

\noindent\textbf{Problem 24: Inscribed Circle via Grand-Yin Right Triangle, 2}

\wenE{A exits the north gate and walks east 50 paces; B exits the east gate and walks north 20 paces. They see each other. Find the city's diameter.}
\daE{The city diameter is 240 paces.}
\shuE{Multiply the grand-yin base by the grand-yin altitude and double it to make the dividend; the sum of grand-yin base and altitude is the divisor; dividing gives the grand-yin inscribed circle diameter; extending by proportion gives the city diameter.}
\shiE{This is finding the inscribed circle via the grand-yin right triangle. Using the grand-yin base and altitude values and reducing them gives the circle's diameter.}

\end{paracol}

\bigskip

\begin{paracol}{2}

\noindent\textbf{第25題 太陰勾股求容圓三}

\wen{甲出西門南行一百步，乙出南門西行四十步，相望見之。問城徑。}
\da{城徑二百四十步。}
\shu{太陰勾股相乘倍之為實，太陰勾股和為法，除之得太陰容圓徑，以比例推之得城徑。}
\shi{此太陰勾股求容圓也。以太陰勾股之數，約之得圓徑。}

\switchcolumn

\noindent\textbf{Problem 25: Inscribed Circle via Grand-Yin Right Triangle, 3}

\wenE{A exits the west gate and walks south 100 paces; B exits the south gate and walks west 40 paces. They see each other. Find the city's diameter.}
\daE{The city diameter is 240 paces.}
\shuE{Multiply the grand-yin base by the grand-yin altitude and double it to make the dividend; the sum of grand-yin base and altitude is the divisor; dividing gives the grand-yin inscribed circle diameter; extending by proportion gives the city diameter.}
\shiE{This is finding the inscribed circle via the grand-yin right triangle. Using the grand-yin base and altitude values and reducing them gives the circle's diameter.}

\end{paracol}

\clearpage

%% ========================================================================
%% Chapter 6: Ming gougu qiu rong yuan
%% ========================================================================
\subsection*{明勾股求容圓 / Finding the Inscribed Circle via the Bright Right Triangle}
\addcontentsline{toc}{subsection}{明勾股求容圓 / Bright Right Triangle}

\begin{paracol}{2}

\noindent\textbf{總論：} 明勾股求容圓者，以明勾、明股之數，用勾股容圓之法，求圓徑也。

\switchcolumn

\noindent\textbf{General Discussion:} To find the inscribed circle via the bright right triangle, use the bright base and bright altitude values, applying the right-triangle inscribed-circle method, to find the circle's diameter.

\end{paracol}

\bigskip

\begin{paracol}{2}

\noindent\textbf{第26題 明勾股求容圓一}

\wen{甲出西門南行八十步，乙出南門東行五十一步，相望見之。問城徑。}
\da{城徑二百四十步。}
\shu{明勾股相乘倍之為實，明勾股和為法，除之得明容圓徑，以比例推之得城徑。}
\shi{此明勾股求容圓也。以明勾股之數，約之得圓徑。}

\switchcolumn

\noindent\textbf{Problem 26: Inscribed Circle via Bright Right Triangle, 1}

\wenE{A exits the west gate and walks south 80 paces; B exits the south gate and walks east 51 paces. They see each other. Find the city's diameter.}
\daE{The city diameter is 240 paces.}
\shuE{Multiply the bright base by the bright altitude and double it to make the dividend; the sum of bright base and altitude is the divisor; dividing gives the bright inscribed circle diameter; extending by proportion gives the city diameter.}
\shiE{This is finding the inscribed circle via the bright right triangle. Using the bright base and altitude values and reducing them gives the circle's diameter.}

\end{paracol}

\bigskip

\begin{paracol}{2}

\noindent\textbf{第27題 明勾股求容圓二}

\wen{甲出北門西行六十步，乙出西門北行三十九步，相望見之。問城徑。}
\da{城徑二百四十步。}
\shu{明勾股相乘倍之為實，明勾股和為法，除之得明容圓徑，以比例推之得城徑。}
\shi{此明勾股求容圓也。以明勾股之數，約之得圓徑。}

\switchcolumn

\noindent\textbf{Problem 27: Inscribed Circle via Bright Right Triangle, 2}

\wenE{A exits the north gate and walks west 60 paces; B exits the west gate and walks north 39 paces. They see each other. Find the city's diameter.}
\daE{The city diameter is 240 paces.}
\shuE{Multiply the bright base by the bright altitude and double it to make the dividend; the sum of bright base and altitude is the divisor; dividing gives the bright inscribed circle diameter; extending by proportion gives the city diameter.}
\shiE{This is finding the inscribed circle via the bright right triangle. Using the bright base and altitude values and reducing them gives the circle's diameter.}

\end{paracol}

\clearpage

%% ========================================================================
%% Chapter 7: Chong gougu qiu rong yuan
%% ========================================================================
\subsection*{重勾股求容圓 / Finding the Inscribed Circle via the Double Right Triangle}
\addcontentsline{toc}{subsection}{重勾股求容圓 / Double Right Triangle}

\begin{paracol}{2}

\noindent\textbf{總論：} 重勾股求容圓者，以重勾、重股之數，用勾股容圓之法，求圓徑也。

\switchcolumn

\noindent\textbf{General Discussion:} To find the inscribed circle via the double right triangle, use the double base and double altitude values, applying the right-triangle inscribed-circle method, to find the circle's diameter.

\end{paracol}

\bigskip

\begin{paracol}{2}

\noindent\textbf{第28題 重勾股求容圓一}

\wen{甲出東門北行三十步，乙出北門西行一十六步，相望見之。問城徑。}
\da{城徑二百四十步。}
\shu{重勾股相乘倍之為實，重勾股和為法，除之得重容圓徑，以比例推之得城徑。}
\shi{此重勾股求容圓也。以重勾股之數，約之得圓徑。}

\switchcolumn

\noindent\textbf{Problem 28: Inscribed Circle via Double Right Triangle, 1}

\wenE{A exits the east gate and walks north 30 paces; B exits the north gate and walks west 16 paces. They see each other. Find the city's diameter.}
\daE{The city diameter is 240 paces.}
\shuE{Multiply the double base by the double altitude and double it to make the dividend; the sum of double base and altitude is the divisor; dividing gives the double inscribed circle diameter; extending by proportion gives the city diameter.}
\shiE{This is finding the inscribed circle via the double right triangle. Using the double base and altitude values and reducing them gives the circle's diameter.}

\end{paracol}

\bigskip

\begin{paracol}{2}

\noindent\textbf{第29題 重勾股求容圓二}

\wen{甲出南門東行二十四步，乙出東門南行十步，相望見之。問城徑。}
\da{城徑二百四十步。}
\shu{重勾股相乘倍之為實，重勾股和為法，除之得重容圓徑，以比例推之得城徑。}
\shi{此重勾股求容圓也。以重勾股之數，約之得圓徑。}

\switchcolumn

\noindent\textbf{Problem 29: Inscribed Circle via Double Right Triangle, 2}

\wenE{A exits the south gate and walks east 24 paces; B exits the east gate and walks south 10 paces. They see each other. Find the city's diameter.}
\daE{The city diameter is 240 paces.}
\shuE{Multiply the double base by the double altitude and double it to make the dividend; the sum of double base and altitude is the divisor; dividing gives the double inscribed circle diameter; extending by proportion gives the city diameter.}
\shiE{This is finding the inscribed circle via the double right triangle. Using the double base and altitude values and reducing them gives the circle's diameter.}

\end{paracol}

\clearpage

%% ========================================================================
%% Reference Table for Volume 1
%% ========================================================================
\subsection*{卷一名數表 / Volume 1 Reference Table}
\addcontentsline{toc}{subsection}{卷一名數表 / Volume 1 Reference Table}

\begin{center}
\textbf{Chinese Terms \& Numerical Values}
\end{center}

\begin{center}
\begin{longtable}{|l|l|l|}
\hline
\textbf{Name} & \textbf{Value} & \textbf{Note} \\
\hline
\endfirsthead
\hline
\textbf{Name} & \textbf{Value} & \textbf{Note} \\
\hline
\endhead
勾股和 (gou gu he) & 736 & Sum of universal base and altitude \\
\hline
勾弦和 (gou xian he) & 800 & Sum of universal base and hypotenuse \\
\hline
股弦和 (gu xian he) & 1024 & Sum of universal altitude and hypotenuse \\
\hline
弦較和 (xian jiao he) & 768 & Sum of hypotenuse and difference \\
\hline
弦較較 (xian jiao jiao) & 288 & Difference of hypotenuse and difference \\
\hline
弦和和 (xian he he) & 1280 & Sum of hypotenuse and sum \\
\hline
弦和較 (xian he jiao) & 224 & Difference of hypotenuse and sum \\
\hline
黃廣弦 (huang guang xian) & 510 & Base 240, altitude 450 \\
\hline
勾股和 (gou gu he) & 690 & Base 240, altitude 450 \\
\hline
勾弦和 (gou xian he) & 750 & Base 240, hypotenuse 510 \\
\hline
較 (jiao) & 210 & Difference of base and altitude \\
\hline
\end{longtable}
\end{center}

\clearpage

%% ========================================================================
%% VOLUME 2 (卷二)
%% ========================================================================
\section{卷二 / Volume 2}
\markboth{卷二 / Volume 2}{卷二 / Volume 2}

\begin{paracol}{2}

\noindent\textbf{卷二}

\smallskip

卷二專論\textbf{兩股求容圓}之法，凡七條。兩股求容圓者，以兩股之數，用通勾股之法，求圓城之徑也。又論\textbf{兩弦求容圓}之法，凡三條。

\switchcolumn

\noindent\textbf{Volume 2}

\smallskip

Volume 2 is devoted to the method of \textbf{finding the inscribed circle using two altitudes} (liang gu qiu rong yuan), comprising seven problems. This method uses two altitude values, applying the universal right triangle method, to find the diameter of the circular city. It also treats the method of \textbf{finding the inscribed circle using two hypotenuses} (liang xian qiu rong yuan), comprising three problems.

\end{paracol}

\bigskip

%% ========================================================================
%% Chapter 1: Liang gu qiu rong yuan
%% ========================================================================
\subsection*{兩股求容圓 / Finding the Inscribed Circle via Two Altitudes}
\addcontentsline{toc}{subsection}{兩股求容圓 / Two Altitudes}

\begin{paracol}{2}

\noindent\textbf{總論：} 兩股求容圓者，以兩股之數相乘倍之為實，以兩股之和為法，除之得圓徑。

\switchcolumn

\noindent\textbf{General Discussion:} To find the inscribed circle via two altitudes, multiply the two altitude values and double it to make the dividend; take the sum of the two altitudes as the divisor; dividing gives the circle's diameter.

\end{paracol}

\bigskip

\begin{paracol}{2}

\noindent\textbf{第30題 兩股求容圓一}

\wen{城南有槐一株，城東有柳一株。甲出北門東行，丙出西門南行。既而甲斜行四百二十五步至槐下，丙斜行五百四十四步至柳下。問城徑。}
\da{城徑二百四十步。}
\shu{二斜行相減餘自之，得一萬四千一百六十一為差冪。甲斜行自之，得一十八萬零六百二十五為冪。以差冪減之，餘為實。以二斜行相減為法，除之得半徑。}
\shi{此以兩股之差求容圓也。甲斜行向西南至槐樹下，底弦也；丙斜行向東北至柳樹下，邊弦也。此以邊弦底弦互測圓徑。}

\switchcolumn

\noindent\textbf{Problem 30: Inscribed Circle via Two Altitudes, 1}

\wenE{South of the city there is a locust tree; east of the city there is a willow tree. A exits the north gate and walks east; C exits the west gate and walks south. Then A walks diagonally 425 paces to reach the locust tree; C walks diagonally 544 paces to reach the willow tree. Find the city's diameter.}
\daE{The city diameter is 240 paces.}
\shuE{Subtract the two diagonal walks; square the remainder to obtain 14,161 as the difference-square. Square A's diagonal walk to obtain 180,625 as the square. Subtract the difference-square from it; the remainder is the dividend. Take the difference of the two diagonal walks as the divisor; dividing gives the radius.}
\shiE{This finds the inscribed circle via the difference of two altitudes. A's diagonal walk southwest to the locust tree is the base-hypotenuse; C's diagonal walk northeast to the willow tree is the edge-hypotenuse. This uses the edge-hypotenuse and base-hypotenuse to mutually measure the circle's diameter.}

\end{paracol}

\bigskip

\begin{paracol}{2}

\noindent\textbf{第31題 兩股求容圓二}

\wen{甲出南門直行一百三十五步而立，乙從城外西北乾隅南行六百步望乙與城相參直。問城徑。}
\da{城徑二百四十步。}
\shu{兩股相乘倍之為實，併兩股為法，除之得圓徑。}
\shi{此兩股求容圓也。以兩股之數，約之得圓徑。}

\switchcolumn

\noindent\textbf{Problem 31: Inscribed Circle via Two Altitudes, 2}

\wenE{A exits the south gate and walks straight 135 paces, then stands still. B walks south 600 paces from the northwest corner (Qian position) outside the city. [A] sees B aligned with the city edge. Find the city's diameter.}
\daE{The city diameter is 240 paces.}
\shuE{Multiply the two altitude values and double it to make the dividend; combine the two altitudes as the divisor; dividing gives the circle's diameter.}
\shiE{This is finding the inscribed circle via two altitudes. Using the two altitude values and reducing them gives the circle's diameter.}

\end{paracol}

\bigskip

\begin{paracol}{2}

\noindent\textbf{第32題 兩股求容圓三}

\wen{甲出東門直行七十二步，乙出北門直行一百二十步，相望見之。問城徑。}
\da{城徑二百四十步。}
\shu{兩股相乘倍之為實，併兩股為法，除之得圓徑。}
\shi{此兩股求容圓也。以兩股之數，約之得圓徑。}

\switchcolumn

\noindent\textbf{Problem 32: Inscribed Circle via Two Altitudes, 3}

\wenE{A exits the east gate and walks straight 72 paces; B exits the north gate and walks straight 120 paces. They see each other. Find the city's diameter.}
\daE{The city diameter is 240 paces.}
\shuE{Multiply the two altitude values and double it to make the dividend; combine the two altitudes as the divisor; dividing gives the circle's diameter.}
\shiE{This is finding the inscribed circle via two altitudes. Using the two altitude values and reducing them gives the circle's diameter.}

\end{paracol}

\bigskip

\begin{paracol}{2}

\noindent\textbf{第33題 兩股求容圓四}

\wen{甲出西門南行一百步，乙出南門東行六十步，相望見之。問城徑。}
\da{城徑二百四十步。}
\shu{兩股相乘倍之為實，併兩股為法，除之得圓徑。}
\shi{此兩股求容圓也。以兩股之數，約之得圓徑。}

\switchcolumn

\noindent\textbf{Problem 33: Inscribed Circle via Two Altitudes, 4}

\wenE{A exits the west gate and walks south 100 paces; B exits the south gate and walks east 60 paces. They see each other. Find the city's diameter.}
\daE{The city diameter is 240 paces.}
\shuE{Multiply the two altitude values and double it to make the dividend; combine the two altitudes as the divisor; dividing gives the circle's diameter.}
\shiE{This is finding the inscribed circle via two altitudes. Using the two altitude values and reducing them gives the circle's diameter.}

\end{paracol}

\bigskip

\begin{paracol}{2}

\noindent\textbf{第34題 兩股求容圓五}

\wen{甲出北門東行二百步，乙出東門南行一百五十步，相望見之。問城徑。}
\da{城徑二百四十步。}
\shu{兩股相乘倍之為實，併兩股為法，除之得圓徑。}
\shi{此兩股求容圓也。以兩股之數，約之得圓徑。}

\switchcolumn

\noindent\textbf{Problem 34: Inscribed Circle via Two Altitudes, 5}

\wenE{A exits the north gate and walks east 200 paces; B exits the east gate and walks south 150 paces. They see each other. Find the city's diameter.}
\daE{The city diameter is 240 paces.}
\shuE{Multiply the two altitude values and double it to make the dividend; combine the two altitudes as the divisor; dividing gives the circle's diameter.}
\shiE{This is finding the inscribed circle via two altitudes. Using the two altitude values and reducing them gives the circle's diameter.}

\end{paracol}

\bigskip

\begin{paracol}{2}

\noindent\textbf{第35題 兩股求容圓六}

\wen{甲出南門直行二百步，乙出東門直行三百步，相望見之。問城徑。}
\da{城徑二百四十步。}
\shu{兩股相乘倍之為實，併兩股為法，除之得圓徑。}
\shi{此兩股求容圓也。以兩股之數，約之得圓徑。}

\switchcolumn

\noindent\textbf{Problem 35: Inscribed Circle via Two Altitudes, 6}

\wenE{A exits the south gate and walks straight 200 paces; B exits the east gate and walks straight 300 paces. They see each other. Find the city's diameter.}
\daE{The city diameter is 240 paces.}
\shuE{Multiply the two altitude values and double it to make the dividend; combine the two altitudes as the divisor; dividing gives the circle's diameter.}
\shiE{This is finding the inscribed circle via two altitudes. Using the two altitude values and reducing them gives the circle's diameter.}

\end{paracol}

\bigskip

\begin{paracol}{2}

\noindent\textbf{第36題 兩股求容圓七}

\wen{甲出西門直行一百八十步，乙出北門直行二百四十步，相望見之。問城徑。}
\da{城徑二百四十步。}
\shu{兩股相乘倍之為實，併兩股為法，除之得圓徑。}
\shi{此兩股求容圓也。以兩股之數，約之得圓徑。}

\switchcolumn

\noindent\textbf{Problem 36: Inscribed Circle via Two Altitudes, 7}

\wenE{A exits the west gate and walks straight 180 paces; B exits the north gate and walks straight 240 paces. They see each other. Find the city's diameter.}
\daE{The city diameter is 240 paces.}
\shuE{Multiply the two altitude values and double it to make the dividend; combine the two altitudes as the divisor; dividing gives the circle's diameter.}
\shiE{This is finding the inscribed circle via two altitudes. Using the two altitude values and reducing them gives the circle's diameter.}

\end{paracol}

\clearpage

%% ========================================================================
%% Chapter 2: Liang xian qiu rong yuan
%% ========================================================================
\subsection*{兩弦求容圓 / Finding the Inscribed Circle via Two Hypotenuses}
\addcontentsline{toc}{subsection}{兩弦求容圓 / Two Hypotenuses}

\begin{paracol}{2}

\noindent\textbf{總論：} 兩弦求容圓者，以兩弦之數，用減從翻法開平方，求圓城之徑也。

\switchcolumn

\noindent\textbf{General Discussion:} To find the inscribed circle via two hypotenuses, use the two hypotenuse values, applying the subtractive-associate inverted method of root extraction, to find the diameter of the circular city.

\end{paracol}

\bigskip

\begin{paracol}{2}

\noindent\textbf{第37題 兩弦求容圓一}

\wen{甲出北門東行，乙出東門南行，相望見之。甲斜行五百步，乙斜行三百步。問城徑。}
\da{城徑二百四十步。}
\shu{兩弦相乘為實，兩弦和為法，除之得圓徑。}
\shi{此兩弦求容圓也。以兩弦之數，約之得圓徑。}

\switchcolumn

\noindent\textbf{Problem 37: Inscribed Circle via Two Hypotenuses, 1}

\wenE{A exits the north gate and walks east; B exits the east gate and walks south. They see each other. A walks diagonally 500 paces; B walks diagonally 300 paces. Find the city's diameter.}
\daE{The city diameter is 240 paces.}
\shuE{Multiply the two hypotenuses to make the dividend; the sum of the two hypotenuses is the divisor; dividing gives the circle's diameter.}
\shiE{This is finding the inscribed circle via two hypotenuses. Using the two hypotenuse values and reducing them gives the circle's diameter.}

\end{paracol}

\bigskip

\begin{paracol}{2}

\noindent\textbf{第38題 兩弦求容圓二}

\wen{甲出南門直行，乙出西門直行，相望見之。甲斜行四百步，乙斜行二百五十步。問城徑。}
\da{城徑二百四十步。}
\shu{兩弦相乘為實，兩弦和為法，除之得圓徑。}
\shi{此兩弦求容圓也。以兩弦之數，約之得圓徑。}

\switchcolumn

\noindent\textbf{Problem 38: Inscribed Circle via Two Hypotenuses, 2}

\wenE{A exits the south gate and walks straight; B exits the west gate and walks straight. They see each other. A walks diagonally 400 paces; B walks diagonally 250 paces. Find the city's diameter.}
\daE{The city diameter is 240 paces.}
\shuE{Multiply the two hypotenuses to make the dividend; the sum of the two hypotenuses is the divisor; dividing gives the circle's diameter.}
\shiE{This is finding the inscribed circle via two hypotenuses. Using the two hypotenuse values and reducing them gives the circle's diameter.}

\end{paracol}

\bigskip

\begin{paracol}{2}

\noindent\textbf{第39題 兩弦求容圓三}

\wen{甲出東門直行，乙出北門直行，相望見之。甲斜行三百六十步，乙斜行二百步。問城徑。}
\da{城徑二百四十步。}
\shu{兩弦相乘為實，兩弦和為法，除之得圓徑。}
\shi{此兩弦求容圓也。以兩弦之數，約之得圓徑。}

\switchcolumn

\noindent\textbf{Problem 39: Inscribed Circle via Two Hypotenuses, 3}

\wenE{A exits the east gate and walks straight; B exits the north gate and walks straight. They see each other. A walks diagonally 360 paces; B walks diagonally 200 paces. Find the city's diameter.}
\daE{The city diameter is 240 paces.}
\shuE{Multiply the two hypotenuses to make the dividend; the sum of the two hypotenuses is the divisor; dividing gives the circle's diameter.}
\shiE{This is finding the inscribed circle via two hypotenuses. Using the two hypotenuse values and reducing them gives the circle's diameter.}

\end{paracol}

\clearpage

%% ========================================================================
%% Chapter 3: Chang duan er gou qiu chengyuanjing
%% ========================================================================
\subsection*{長短二勾求城徑與通股小差股同法 / Finding City Diameter via Long and Short Bases}
\addcontentsline{toc}{subsection}{長短二勾求城徑 / Long and Short Bases}

\begin{paracol}{2}

\noindent\textbf{總論：} 長短二勾求城徑者，以長勾、短勾之數，用通股小差股之法，求圓城之徑也。

\switchcolumn

\noindent\textbf{General Discussion:} To find the city diameter via long and short bases, use the long-base and short-base values, applying the method of universal-altitude small-difference altitude, to find the diameter of the circular city.

\end{paracol}

\bigskip

\begin{paracol}{2}

\noindent\textbf{第40題 長短二勾求城徑一}

\wen{甲出北門東行二百步，乙出東門南行一百五十步，相望見之。問城徑。}
\da{城徑二百四十步。}
\shu{二行相乘倍之為實，併二行為法，除之得城徑。}
\shi{此長短二勾求城徑也。以長短二勾之數，約之得圓徑。}

\switchcolumn

\noindent\textbf{Problem 40: City Diameter via Long and Short Bases, 1}

\wenE{A exits the north gate and walks east 200 paces; B exits the east gate and walks south 150 paces. They see each other. Find the city's diameter.}
\daE{The city diameter is 240 paces.}
\shuE{Multiply the two walks and double it to make the dividend; combine the two walks as the divisor; dividing gives the city diameter.}
\shiE{This is finding the city diameter via long and short bases. Using the long and short base values and reducing them gives the circle's diameter.}

\end{paracol}

\bigskip

\begin{paracol}{2}

\noindent\textbf{第41題 長短二勾求城徑二}

\wen{甲出南門直行一百八十步，乙出西門直行一百二十步，相望見之。問城徑。}
\da{城徑二百四十步。}
\shu{二行相乘倍之為實，併二行為法，除之得城徑。}
\shi{此長短二勾求城徑也。以長短二勾之數，約之得圓徑。}

\switchcolumn

\noindent\textbf{Problem 41: City Diameter via Long and Short Bases, 2}

\wenE{A exits the south gate and walks straight 180 paces; B exits the west gate and walks straight 120 paces. They see each other. Find the city's diameter.}
\daE{The city diameter is 240 paces.}
\shuE{Multiply the two walks and double it to make the dividend; combine the two walks as the divisor; dividing gives the city diameter.}
\shiE{This is finding the city diameter via long and short bases. Using the long and short base values and reducing them gives the circle's diameter.}

\end{paracol}

\clearpage

%% ========================================================================
%% Chapter 4: Jian cong fanfa kaipingfang
%% ========================================================================
\subsection*{減從翻法開平方 / Subtractive-Associate Inverted Root Extraction}
\addcontentsline{toc}{subsection}{減從翻法開平方 / Subtractive-Associate Inverted Root Extraction}

\begin{paracol}{2}

\noindent\textbf{總論：} 減從翻法開平方者，以實較大於法，用減從之法，翻轉開平方，得圓城之徑也。

\switchcolumn

\noindent\textbf{General Discussion:} The subtractive-associate inverted method of root extraction is used when the dividend exceeds the divisor. Using the subtractive-associate method, one inverts and opens the square to obtain the diameter of the circular city.

\end{paracol}

\bigskip

\begin{paracol}{2}

\noindent\textbf{第42題 減從翻法開平方一}

\wen{餘實五百五十萬，倍隅法得二十五萬為廉法。約次商得二十，置一於左次為上法，置一乘隅算得二萬五千，併廉法共二十七萬五千為下法，與上法相乘除實盡。後如此類者倣此。}
\da{得半徑一百二十步。}
\shu{以減從翻法開平方，置一於左上為上法，倍隅法為廉法，次商自之為隅法，併廉隅共為下法，與上法相乘除實。}
\shi{此減從翻法開平方也。以減從翻法開平方，得圓徑。}

\switchcolumn

\noindent\textbf{Problem 42: Subtractive-Associate Inverted Root Extraction, 1}

\wenE{The remaining dividend is 5,500,000; double the corner-method to obtain 250,000 as the edge-method. Estimating the next quotient as 20, place one on the left as the next upper-method; place one and multiply by the corner-count to obtain 25,000; combine with the edge-method for a total of 275,000 as the lower-method; multiply by the upper-method to exhaust the dividend. All later cases of this type follow this pattern.}
\daE{The radius obtained is 120 paces.}
\shuE{Using the subtractive-associate inverted method of root extraction: place one in the upper left as the upper-method; double the corner-method as the edge-method; the next quotient squared is the corner-method; combine edge and corner as the lower-method; multiply by the upper-method to divide the dividend.}
\shiE{This is the subtractive-associate inverted method of root extraction. Using this method to open the square gives the circle's diameter.}

\end{paracol}

\bigskip

\begin{paracol}{2}

\noindent\textbf{第43題 減從翻法開平方二}

\wen{二十與上次法相乘，得四千四百為益實。添入餘積，共五千六百為實。置一併廉法，從方共二百八十為下法，與上次法相乘除實盡。}
\da{得半徑一百二十步。}
\shu{以減從翻法開平方，從方添入益實，併廉法為下法，與上法相乘除實。}
\shi{此減從翻法開平方也。以減從翻法開平方，得圓徑。}

\switchcolumn

\noindent\textbf{Problem 43: Subtractive-Associate Inverted Root Extraction, 2}

\wenE{Twenty multiplied by the previous upper-method gives 4,400 as the augmenting dividend. Add it to the remaining accumulation; the total of 5,600 is the dividend. Place one combined with the edge-method; the associate-square totals 280 as the lower-method; multiply by the previous upper-method to exhaust the dividend.}
\daE{The radius obtained is 120 paces.}
\shuE{Using the subtractive-associate inverted method of root extraction: add the associate-square into the augmenting dividend; combine the edge-method as the lower-method; multiply by the upper-method to divide the dividend.}
\shiE{This is the subtractive-associate inverted method of root extraction. Using this method to open the square gives the circle's diameter.}

\end{paracol}

\clearpage

%% ========================================================================
%% Chapter 5: Gou gu xiang cheng beizhi fa
%% ========================================================================
\subsection*{勾股相乘倍之為實開平方 / Doubling the Product of Base and Altitude as Dividend}
\addcontentsline{toc}{subsection}{勾股相乘倍之為實開平方 / Doubled Product Root Extraction}

\begin{paracol}{2}

\noindent\textbf{總論：} 勾股相乘倍之為實者，以勾股相乘之數倍之為實，開平方得弦，再以弦求容圓也。

\switchcolumn

\noindent\textbf{General Discussion:} Doubling the product of base and altitude as dividend: take the product of base and altitude, double it as the dividend, open the square root to obtain the hypotenuse, then use the hypotenuse to find the inscribed circle.

\end{paracol}

\bigskip

\begin{paracol}{2}

\noindent\textbf{第44題 勾股相乘倍之開平方一}

\wen{甲出南門東行一百二十步，乙出東門南行八十步，相望見之。問城徑。}
\da{城徑二百四十步。}
\shu{勾股相乘倍之為實，開平方得弦，併勾股為弦和，弦和較即容圓徑。}
\shi{此勾股相乘倍之為實開平方也。以勾股相乘倍之為實，開平方得弦，再以弦求容圓。}

\switchcolumn

\noindent\textbf{Problem 44: Doubled Product Root Extraction, 1}

\wenE{A exits the south gate and walks east 120 paces; B exits the east gate and walks south 80 paces. They see each other. Find the city's diameter.}
\daE{The city diameter is 240 paces.}
\shuE{Multiply base and altitude, double it as the dividend; open the square root to obtain the hypotenuse; combine base and altitude as the hypotenuse-sum; the difference between hypotenuse and sum is the inscribed circle's diameter.}
\shiE{This is the method of doubling the product of base and altitude as dividend. Doubling the product as the dividend and opening the square gives the hypotenuse; then use the hypotenuse to find the inscribed circle.}

\end{paracol}

\bigskip

\begin{paracol}{2}

\noindent\textbf{第45題 勾股相乘倍之開平方二}

\wen{甲出北門東行一百五十步，乙出東門北行一百步，相望見之。問城徑。}
\da{城徑二百四十步。}
\shu{勾股相乘倍之為實，開平方得弦，併勾股為弦和，弦和較即容圓徑。}
\shi{此勾股相乘倍之為實開平方也。以勾股相乘倍之為實，開平方得弦，再以弦求容圓。}

\switchcolumn

\noindent\textbf{Problem 45: Doubled Product Root Extraction, 2}

\wenE{A exits the north gate and walks east 150 paces; B exits the east gate and walks north 100 paces. They see each other. Find the city's diameter.}
\daE{The city diameter is 240 paces.}
\shuE{Multiply base and altitude, double it as the dividend; open the square root to obtain the hypotenuse; combine base and altitude as the hypotenuse-sum; the difference between hypotenuse and sum is the inscribed circle's diameter.}
\shiE{This is the method of doubling the product of base and altitude as dividend. Doubling the product as the dividend and opening the square gives the hypotenuse; then use the hypotenuse to find the inscribed circle.}

\end{paracol}

\clearpage

%% ========================================================================
%% Chapter 6: Xian he jiao qiu yuanjing
%% ========================================================================
\subsection*{弦和較求圓徑 / Finding Circle Diameter via Hypotenuse-Sum Difference}
\addcontentsline{toc}{subsection}{弦和較求圓徑 / Hypotenuse-Sum Difference}

\begin{paracol}{2}

\noindent\textbf{總論：} 弦和較求圓徑者，以弦與勾股和之差，即容圓徑也。

\switchcolumn

\noindent\textbf{General Discussion:} To find the circle diameter via the hypotenuse-sum difference, the difference between the hypotenuse and the sum of base and altitude is the inscribed circle's diameter.

\end{paracol}

\bigskip

\begin{paracol}{2}

\noindent\textbf{第46題 弦和較求圓徑一}

\wen{甲出南門直行，乙出東門直行，相望見之。問以弦和較求圓徑。}
\da{城徑二百四十步。}
\shu{弦和較即容圓徑。弦者勾股弦也，和者勾股和也。弦減勾股和之較即圓徑。}
\shi{此弦和較求圓徑也。弦和較者，通弦減通勾股和之較也，即容圓徑。}

\switchcolumn

\noindent\textbf{Problem 46: Circle Diameter via Hypotenuse-Sum Difference, 1}

\wenE{A exits the south gate and walks straight; B exits the east gate and walks straight. They see each other. Find the circle diameter using the hypotenuse-sum difference.}
\daE{The city diameter is 240 paces.}
\shuE{The hypotenuse-sum difference is the inscribed circle's diameter. The hypotenuse is the right-triangle hypotenuse; the sum is the base-altitude sum. The difference of the hypotenuse subtracted from the base-altitude sum is the circle's diameter.}
\shiE{This is finding the circle diameter via the hypotenuse-sum difference. The hypotenuse-sum difference is the difference of the universal hypotenuse minus the universal base-altitude sum, which equals the inscribed circle's diameter.}

\end{paracol}

\bigskip

\begin{paracol}{2}

\noindent\textbf{第47題 弦和較求圓徑二}

\wen{甲出北門直行，乙出西門直行，相望見之。問以弦和較求圓徑。}
\da{城徑二百四十步。}
\shu{弦和較即容圓徑。弦者勾股弦也，和者勾股和也。弦減勾股和之較即圓徑。}
\shi{此弦和較求圓徑也。以弦和較之數，即容圓徑。}

\switchcolumn

\noindent\textbf{Problem 47: Circle Diameter via Hypotenuse-Sum Difference, 2}

\wenE{A exits the north gate and walks straight; B exits the west gate and walks straight. They see each other. Find the circle diameter using the hypotenuse-sum difference.}
\daE{The city diameter is 240 paces.}
\shuE{The hypotenuse-sum difference is the inscribed circle's diameter. The hypotenuse is the right-triangle hypotenuse; the sum is the base-altitude sum. The difference of the hypotenuse subtracted from the base-altitude sum is the circle's diameter.}
\shiE{This is finding the circle diameter via the hypotenuse-sum difference. The hypotenuse-sum difference value is the inscribed circle's diameter.}

\end{paracol}

\clearpage

%% ========================================================================
%% Reference Table for Volume 2
%% ========================================================================
\subsection*{卷二名數表 / Volume 2 Reference Table}
\addcontentsline{toc}{subsection}{卷二名數表 / Volume 2 Reference Table}

\begin{center}
\textbf{Chinese Terms \& Numerical Values}
\end{center}

\begin{center}
\begin{longtable}{|l|l|l|}
\hline
\textbf{Name} & \textbf{Value} & \textbf{Note} \\
\hline
\endfirsthead
\hline
\textbf{Name} & \textbf{Value} & \textbf{Note} \\
\hline
\endhead
通股 (tong gu) & 480 & Universal altitude \\
\hline
邊股 (bian gu) & 240 & Edge altitude \\
\hline
底股 (di gu) & 360 & Base altitude \\
\hline
黃廣股 (huang guang gu) & 450 & Yellow-broad altitude \\
\hline
黃長股 (huang chang gu) & 120 & Yellow-long altitude \\
\hline
高股 (gao gu) & 200 & High altitude \\
\hline
平股 (ping gu) & 80 & Level altitude \\
\hline
大差股 (da cha gu) & 160 & Great-difference altitude \\
\hline
小差股 (xiao cha gu) & 60 & Small-difference altitude \\
\hline
皇極股 (huang ji gu) & 240 & Supreme-pole altitude \\
\hline
太虛股 (tai xu gu) & 30 & Grand-void altitude \\
\hline
明股 (ming gu) & 80 & Bright altitude \\
\hline
重股 (chong gu) & 16 & Double altitude \\
\hline
\end{longtable}
\end{center}

\clearpage

%% ========================================================================
%% VOLUME 3 (卷三)
%% ========================================================================
\section{卷三 / Volume 3}
\markboth{卷三 / Volume 3}{卷三 / Volume 3}

\begin{paracol}{2}

\noindent\textbf{卷三}

\smallskip

卷三專論\textbf{帶從開平方法}及\textbf{三乘方}之法，凡十餘條。帶從開平方法者，以實較大於法，用帶從之法，開平方得圓城之徑也。又論減從翻法、益積、減積諸法，皆開方之變術也。

\switchcolumn

\noindent\textbf{Volume 3}

\smallskip

Volume 3 is devoted to the method of \textbf{accompanying-associate root extraction} (dai cong kai ping fang fa) and the method of \textbf{triple multiplication} (san cheng fang), comprising more than ten problems. The accompanying-associate root extraction method is used when the dividend exceeds the divisor; using the accompanying-associate method, one opens the square to obtain the diameter of the circular city. It also treats the subtractive-associate inverted method, additive accumulation, and subtractive accumulation---all variants of root extraction.

\end{paracol}

\bigskip

%% ========================================================================
%% Chapter 1: Dai cong kaipingfang fa
%% ========================================================================
\subsection*{帶從開平方法 / Accompanying-Associate Root Extraction}
\addcontentsline{toc}{subsection}{帶從開平方法 / Accompanying-Associate Root Extraction}

\begin{paracol}{2}

\noindent\textbf{總論：} 帶從開平方法者，以帶從之實，用開平方之法，得圓城之徑也。所謂帶從者，實內帶有從方之數也。

\switchcolumn

\noindent\textbf{General Discussion:} The accompanying-associate method of root extraction uses a dividend that carries an accompanying-associate value, applying the method of root extraction to obtain the diameter of the circular city. The term ``accompanying-associate'' means the dividend contains an accompanying-square value.

\end{paracol}

\bigskip

\begin{paracol}{2}

\noindent\textbf{第48題 帶從開平方法一}

\wen{乙出南門直行一百三十五步，甲出北門東行二百步見之。問以帶從開平方求城徑。}
\da{城徑二百四十步。}
\shu{為實以南行三百六十為從方，作帶從開平方法除之得全徑。}
\shi{此帶從開平方法也。以從方之數，開平方得圓徑。}

\switchcolumn

\noindent\textbf{Problem 48: Accompanying-Associate Root Extraction, 1}

\wenE{B exits the south gate and walks straight 135 paces; A exits the north gate and walks east 200 paces and sees B. Find the city diameter using accompanying-associate root extraction.}
\daE{The city diameter is 240 paces.}
\shuE{Make it the dividend; take the southward walk of 360 as the accompanying-square; perform accompanying-associate root extraction, dividing to obtain the full diameter.}
\shiE{This is the accompanying-associate method of root extraction. Using the accompanying-square value and opening the square gives the circle's diameter.}

\end{paracol}

\bigskip

\begin{paracol}{2}

\noindent\textbf{第49題 帶從開平方法二}

\wen{甲出北門東行一百二十步，乙出東門南行一百步，相望見之。問以帶從開平方求城徑。}
\da{城徑二百四十步。}
\shu{為實以兩行之和為從方，作帶從開平方法除之得全徑。}
\shi{此帶從開平方法也。以從方之數，開平方得圓徑。}

\switchcolumn

\noindent\textbf{Problem 49: Accompanying-Associate Root Extraction, 2}

\wenE{A exits the north gate and walks east 120 paces; B exits the east gate and walks south 100 paces. They see each other. Find the city diameter using accompanying-associate root extraction.}
\daE{The city diameter is 240 paces.}
\shuE{Make it the dividend; take the sum of the two walks as the accompanying-square; perform accompanying-associate root extraction, dividing to obtain the full diameter.}
\shiE{This is the accompanying-associate method of root extraction. Using the accompanying-square value and opening the square gives the circle's diameter.}

\end{paracol}

\bigskip

\begin{paracol}{2}

\noindent\textbf{第50題 帶從開平方法三}

\wen{甲出西門南行二百步，乙出南門西行一百五十步，相望見之。問以帶從開平方求城徑。}
\da{城徑二百四十步。}
\shu{為實以兩行之和為從方，作帶從開平方法除之得全徑。}
\shi{此帶從開平方法也。以從方之數，開平方得圓徑。}

\switchcolumn

\noindent\textbf{Problem 50: Accompanying-Associate Root Extraction, 3}

\wenE{A exits the west gate and walks south 200 paces; B exits the south gate and walks west 150 paces. They see each other. Find the city diameter using accompanying-associate root extraction.}
\daE{The city diameter is 240 paces.}
\shuE{Make it the dividend; take the sum of the two walks as the accompanying-square; perform accompanying-associate root extraction, dividing to obtain the full diameter.}
\shiE{This is the accompanying-associate method of root extraction. Using the accompanying-square value and opening the square gives the circle's diameter.}

\end{paracol}

\bigskip

\begin{paracol}{2}

\noindent\textbf{第51題 帶從開平方法四}

\wen{甲出南門直行二百五十步，乙出東門直行二百步，相望見之。問以帶從開平方求城徑。}
\da{城徑二百四十步。}
\shu{為實以兩行之和為從方，作帶從開平方法除之得全徑。}
\shi{此帶從開平方法也。以從方之數，開平方得圓徑。}

\switchcolumn

\noindent\textbf{Problem 51: Accompanying-Associate Root Extraction, 4}

\wenE{A exits the south gate and walks straight 250 paces; B exits the east gate and walks straight 200 paces. They see each other. Find the city diameter using accompanying-associate root extraction.}
\daE{The city diameter is 240 paces.}
\shuE{Make it the dividend; take the sum of the two walks as the accompanying-square; perform accompanying-associate root extraction, dividing to obtain the full diameter.}
\shiE{This is the accompanying-associate method of root extraction. Using the accompanying-square value and opening the square gives the circle's diameter.}

\end{paracol}

\clearpage

%% ========================================================================
%% Chapter 2: Jian cong fanfa kaipingfang
%% ========================================================================
\subsection*{減從翻法開平方 / Subtractive-Associate Inverted Root Extraction}
\addcontentsline{toc}{subsection}{減從翻法開平方 (Vol.~3) / Subtractive-Associate Inverted Root Extraction}

\begin{paracol}{2}

\noindent\textbf{總論：} 減從翻法開平方者，以實內減去從方之數，翻轉開平方，得圓城之徑也。凡言減從翻法開平方者俱倣此。

\switchcolumn

\noindent\textbf{General Discussion:} The subtractive-associate inverted method of root extraction subtracts the accompanying-square value from within the dividend, then inverts and opens the square to obtain the diameter of the circular city. All cases mentioning this method follow this pattern.

\end{paracol}

\bigskip

\begin{paracol}{2}

\noindent\textbf{第52題 減從翻法開平方一}

\wen{十餘一千五百二十為負積。又以初商一百反減餘，從四十四餘五十六為負。從次商二十併負，從共七十六為下法，亦通後凡言減從翻法開平方者俱倣此。}
\da{得半徑一百二十步。}
\shu{以減從翻法開平方曰：初商一百，置一於左上為上法。置一乘隅法，得四為隅法，作負隅開平方。}
\shi{此減從翻法開平方也。以減從翻法開平方，得圓徑。}

\switchcolumn

\noindent\textbf{Problem 52: Subtractive-Associate Inverted Root Extraction, 1}

\wenE{Ten plus 1,520 makes the negative accumulation. Also, subtracting the first quotient of 100 from the remainder, the associate of 44 leaves 56 as negative. The next quotient of 20 combined with the negative associate makes 76 total as the lower-method. This pattern applies to all subsequent cases of subtractive-associate inverted root extraction.}
\daE{The radius obtained is 120 paces.}
\shuE{Using the subtractive-associate inverted method of root extraction: first quotient 100; place one in the upper left as the upper-method. Place one and multiply by the corner-method, obtaining 4 as the corner-method; make it a negative corner and open the square.}
\shiE{This is the subtractive-associate inverted method of root extraction. Using this method to open the square gives the circle's diameter.}

\end{paracol}

\bigskip

\begin{paracol}{2}

\noindent\textbf{第53題 減從翻法開平方二}

\wen{餘實五百五十萬，倍隅法得二十五萬為廉法。約次商得二十，置一於左次為上法，置一乘隅得二萬五千，併廉法共二十七萬五千為下法，與上法相乘除實。盡後如此類者倣此。}
\da{得半徑一百二十步。}
\shu{以減從翻法開平方，置一於左上為上法，倍隅法為廉法，次商自之為隅法，併廉隅共為下法，與上法相乘除實。}
\shi{此減從翻法開平方也。以減從翻法開平方，得圓徑。}

\switchcolumn

\noindent\textbf{Problem 53: Subtractive-Associate Inverted Root Extraction, 2}

\wenE{The remaining dividend is 5,500,000; double the corner-method to obtain 250,000 as the edge-method. Estimating the next quotient as 20, place one on the left as the next upper-method; place one and multiply by the corner to obtain 25,000; combine with the edge-method for a total of 275,000 as the lower-method; multiply by the upper-method to divide the dividend. All later cases of this type follow this pattern.}
\daE{The radius obtained is 120 paces.}
\shuE{Using the subtractive-associate inverted method of root extraction: place one in the upper left as the upper-method; double the corner-method as the edge-method; the next quotient squared is the corner-method; combine edge and corner as the lower-method; multiply by the upper-method to divide the dividend.}
\shiE{This is the subtractive-associate inverted method of root extraction. Using this method to open the square gives the circle's diameter.}

\end{paracol}

\bigskip

\begin{paracol}{2}

\noindent\textbf{第54題 減從翻法開平方三}

\wen{二十與上次法相乘，得四千四百為益實。添入餘積共五千六百為實。置一併廉法從方共二百八十為下法，與上次法相乘除實盡。}
\da{得半徑一百二十步。}
\shu{以減從翻法開平方，從方添入益實，併廉法為下法，與上法相乘除實。}
\shi{此減從翻法開平方也。以減從翻法開平方，得圓徑。}

\switchcolumn

\noindent\textbf{Problem 54: Subtractive-Associate Inverted Root Extraction, 3}

\wenE{Twenty multiplied by the previous upper-method gives 4,400 as the augmenting dividend. Add it to the remaining accumulation; the total of 5,600 is the dividend. Place one combined with the edge-method; the associate-square totals 280 as the lower-method; multiply by the previous upper-method to exhaust the dividend.}
\daE{The radius obtained is 120 paces.}
\shuE{Using the subtractive-associate inverted method of root extraction: add the associate-square into the augmenting dividend; combine the edge-method as the lower-method; multiply by the upper-method to divide the dividend.}
\shiE{This is the subtractive-associate inverted method of root extraction. Using this method to open the square gives the circle's diameter.}

\end{paracol}

\clearpage

%% ========================================================================
%% Chapter 3: San cheng fang fa
%% ========================================================================
\subsection*{三乘方法 / Triple Multiplication Method}
\addcontentsline{toc}{subsection}{三乘方法 / Triple Multiplication Method}

\begin{paracol}{2}

\noindent\textbf{總論：} 三乘方法者，以三次方程之法，求圓城之徑也。所謂三乘者，立方以上之開方也。

\switchcolumn

\noindent\textbf{General Discussion:} The triple multiplication method uses cubic equations to find the diameter of the circular city. The term ``triple multiplication'' refers to root extraction beyond the cube root.

\end{paracol}

\bigskip

\begin{paracol}{2}

\noindent\textbf{第55題 三乘方法一}

\wen{置一自之以乘，從二廉，置一自乘再乘為隅法。併從方廉隅共七千九十八萬一千三百四，併從方廉隅共七千九十八萬一千三百四十為下法，與上法相乘除實。}
\da{得半徑一百二十步。}
\shu{以三乘方法開之，置一自之為上法，以次商乘廉法為從方，再以次商自之為隅法，併從方廉隅為下法，與上法相乘除實。}
\shi{此三乘方法也。以三乘方法開之，得圓徑。}

\switchcolumn

\noindent\textbf{Problem 55: Triple Multiplication Method, 1}

\wenE{Place one, self-multiply it to multiply the two accompanying edges; place one, self-multiply and multiply again as the corner-method. Combining associate-square, edge, and corner makes 70,981,304; combining associate-square, edge, and corner makes 70,981,340 as the lower-method; multiply by the upper-method to divide the dividend.}
\daE{The radius obtained is 120 paces.}
\shuE{Using the triple multiplication method: place one, self-multiply as the upper-method; multiply the next quotient by the edge-method as the associate-square; then square the next quotient as the corner-method; combine associate-square, edge, and corner as the lower-method; multiply by the upper-method to divide the dividend.}
\shiE{This is the triple multiplication method. Using this method gives the circle's diameter.}

\end{paracol}

\bigskip

\begin{paracol}{2}

\noindent\textbf{第56題 三乘方法二}

\wen{置一自之以乘從二廉，置一自乘再乘為隅法。併從方廉隅共七千九十八萬一千三百四十為下法，與上法相乘除實盡。}
\da{得半徑一百二十步。}
\shu{以三乘方法開之，置一自之為上法，以次商乘廉法為從方，再以次商自之為隅法，併從方廉隅為下法，與上法相乘除實。}
\shi{此三乘方法也。以三乘方法開之，得圓徑。}

\switchcolumn

\noindent\textbf{Problem 56: Triple Multiplication Method, 2}

\wenE{Place one, self-multiply it to multiply the two accompanying edges; place one, self-multiply and multiply again as the corner-method. Combining associate-square, edge, and corner makes 70,981,340 as the lower-method; multiply by the upper-method to exhaust the dividend.}
\daE{The radius obtained is 120 paces.}
\shuE{Using the triple multiplication method: place one, self-multiply as the upper-method; multiply the next quotient by the edge-method as the associate-square; then square the next quotient as the corner-method; combine associate-square, edge, and corner as the lower-method; multiply by the upper-method to divide the dividend.}
\shiE{This is the triple multiplication method. Using this method gives the circle's diameter.}

\end{paracol}

\clearpage

%% ========================================================================
%% Chapter 4: Yi ji jian ji fa
%% ========================================================================
\subsection*{益積減積法 / Additive and Subtractive Accumulation Methods}
\addcontentsline{toc}{subsection}{益積減積法 / Additive and Subtractive Accumulation}

\begin{paracol}{2}

\noindent\textbf{總論：} 益積減積法者，以益積或減積之法，開方得圓城之徑也。益積者添入積數，減積者減去積數。

\switchcolumn

\noindent\textbf{General Discussion:} The additive and subtractive accumulation methods use either adding to or subtracting from the accumulated value, then performing root extraction to obtain the diameter of the circular city. Additive accumulation adds to the accumulated number; subtractive accumulation subtracts from it.

\end{paracol}

\bigskip

\begin{paracol}{2}

\noindent\textbf{第57題 益積減積法一}

\wen{二十與上次法相乘，得四千四百為益實。添入餘積，共五千六百為實。置一併廉法，從方共二百八十為下法，與上次法相乘除實盡。}
\da{得半徑一百二十步。}
\shu{以益積法添入益實，併廉法為下法，與上法相乘除實。}
\shi{此益積法也。以益積之數，開平方得圓徑。}

\switchcolumn

\noindent\textbf{Problem 57: Additive and Subtractive Accumulation, 1}

\wenE{Twenty multiplied by the previous upper-method gives 4,400 as the augmenting dividend. Add it to the remaining accumulation; the total of 5,600 is the dividend. Place one combined with the edge-method; the associate-square totals 280 as the lower-method; multiply by the previous upper-method to exhaust the dividend.}
\daE{The radius obtained is 120 paces.}
\shuE{Using the additive accumulation method: add into the augmenting dividend; combine the edge-method as the lower-method; multiply by the upper-method to divide the dividend.}
\shiE{This is the additive accumulation method. Using the additive accumulation value and opening the square gives the circle's diameter.}

\end{paracol}

\bigskip

\begin{paracol}{2}

\noindent\textbf{第58題 益積減積法二}

\wen{餘實五百五十萬，倍隅法得二十五萬為廉法。約次商得二十，置一於左次為上法，置一乘隅得二萬五千，併廉法共二十七萬五千為下法，與上法相乘除實。}
\da{得半徑一百二十步。}
\shu{以減積法減去積數，併廉法為下法，與上法相乘除實。}
\shi{此減積法也。以減積之數，開平方得圓徑。}

\switchcolumn

\noindent\textbf{Problem 58: Additive and Subtractive Accumulation, 2}

\wenE{The remaining dividend is 5,500,000; double the corner-method to obtain 250,000 as the edge-method. Estimating the next quotient as 20, place one on the left as the next upper-method; place one and multiply by the corner to obtain 25,000; combine with the edge-method for a total of 275,000 as the lower-method; multiply by the upper-method to divide the dividend.}
\daE{The radius obtained is 120 paces.}
\shuE{Using the subtractive accumulation method: subtract from the accumulated number; combine the edge-method as the lower-method; multiply by the upper-method to divide the dividend.}
\shiE{This is the subtractive accumulation method. Using the subtractive accumulation value and opening the square gives the circle's diameter.}

\end{paracol}

\clearpage

%% ========================================================================
%% Chapter 5: Cong fang cong lian yu fa
%% ========================================================================
\subsection*{從方從廉隅法 / Associate-Square, Associate-Edge, and Corner Methods}
\addcontentsline{toc}{subsection}{從方從廉隅法 / Associate-Square, Edge, and Corner}

\begin{paracol}{2}

\noindent\textbf{總論：} 從方從廉隅法者，以從方、從廉、隅法三者之關係，開方得圓城之徑也。

\switchcolumn

\noindent\textbf{General Discussion:} The associate-square, associate-edge, and corner method uses the relationships among the three quantities---associate-square, associate-edge, and corner-method---to perform root extraction and obtain the diameter of the circular city.

\end{paracol}

\bigskip

\begin{paracol}{2}

\noindent\textbf{第59題 從方從廉隅法一}

\wen{置一自之以乘從二廉，置一自乘再乘為隅法。併從方廉隅共七千九十八萬一千三百四十為下法，與上法相乘除實盡。}
\da{得半徑一百二十步。}
\shu{以從方從廉隅法開之，置一自之為上法，以次商乘廉法為從方，再以次商自之為隅法，併從方廉隅為下法，與上法相乘除實。}
\shi{此從方從廉隅法也。以從方從廉隅法開之，得圓徑。}

\switchcolumn

\noindent\textbf{Problem 59: Associate-Square, Edge, and Corner Method, 1}

\wenE{Place one, self-multiply it to multiply the two accompanying edges; place one, self-multiply and multiply again as the corner-method. Combining associate-square, edge, and corner makes 70,981,340 as the lower-method; multiply by the upper-method to exhaust the dividend.}
\daE{The radius obtained is 120 paces.}
\shuE{Using the associate-square, associate-edge, and corner method: place one, self-multiply as the upper-method; multiply the next quotient by the edge-method as the associate-square; then square the next quotient as the corner-method; combine associate-square, edge, and corner as the lower-method; multiply by the upper-method to divide the dividend.}
\shiE{This is the associate-square, associate-edge, and corner method. Using this method gives the circle's diameter.}

\end{paracol}

\bigskip

\begin{paracol}{2}

\noindent\textbf{第60題 從方從廉隅法二}

\wen{置一自之為上法，置一乘隅得二萬五千，併廉法共二十七萬五千為下法，與上法相乘除實。}
\da{得半徑一百二十步。}
\shu{以從方從廉隅法開之，置一自之為上法，以次商乘廉法為從方，再以次商自之為隅法，併從方廉隅為下法，與上法相乘除實。}
\shi{此從方從廉隅法也。以從方從廉隅法開之，得圓徑。}

\switchcolumn

\noindent\textbf{Problem 60: Associate-Square, Edge, and Corner Method, 2}

\wenE{Place one, self-multiply as the upper-method; place one and multiply by the corner to obtain 25,000; combine with the edge-method for a total of 275,000 as the lower-method; multiply by the upper-method to divide the dividend.}
\daE{The radius obtained is 120 paces.}
\shuE{Using the associate-square, associate-edge, and corner method: place one, self-multiply as the upper-method; multiply the next quotient by the edge-method as the associate-square; then square the next quotient as the corner-method; combine associate-square, edge, and corner as the lower-method; multiply by the upper-method to divide the dividend.}
\shiE{This is the associate-square, associate-edge, and corner method. Using this method gives the circle's diameter.}

\end{paracol}

\clearpage

%% ========================================================================
%% Chapter 6: Bu ji di gou bian gu tiao
%% ========================================================================
\subsection*{不及底勾邊股條 / Edge-Leg and Base-Leg Deficiency Method}
\addcontentsline{toc}{subsection}{不及底勾邊股條 / Edge-Leg and Base-Leg Deficiency}

\begin{paracol}{2}

\noindent\textbf{總論：} 不及底勾邊股條者，以底勾、邊股不及之數，用測望之法，求圓城之徑也。

\switchcolumn

\noindent\textbf{General Discussion:} The edge-leg and base-leg deficiency method uses the amounts by which the base-leg and edge-leg fall short, applying the observation method, to find the diameter of the circular city.

\end{paracol}

\bigskip

\begin{paracol}{2}

\noindent\textbf{第61題 不及底勾邊股條一}

\wen{出西門南行二百二十五步，有塔出北門東行六十四步望塔正居城之半。問城徑。}
\da{城徑二百四十步。}
\shu{此以不及底勾與不及邊股測望也。南行二百二十五步與高股同，即半徑為勾之股。東行六十四步與平勾同，即半徑為股之勾也。當以平勾高股立法。}
\shi{此以不及底勾邊股條求城徑也。以底勾邊股之不及數，測望得圓徑。}

\switchcolumn

\noindent\textbf{Problem 61: Edge-Leg and Base-Leg Deficiency, 1}

\wenE{Exiting the west gate and walking south 225 paces, there is a pagoda. Exiting the north gate and walking east 64 paces, one sees the pagoda exactly at the center of the city. Find the city's diameter.}
\daE{The city diameter is 240 paces.}
\shuE{This uses the deficiency of base-leg and deficiency of edge-leg for observation. The southward walk of 225 paces equals the high-altitude; that is, the radius serves as the altitude to the base. The eastward walk of 64 paces equals the level-base; that is, the radius serves as the base to the altitude. One should establish the method using level-base and high-altitude.}
\shiE{This finds the city diameter via the edge-leg and base-leg deficiency method. Using the deficient amounts of base-leg and edge-leg, observation yields the circle's diameter.}

\end{paracol}

\bigskip

\begin{paracol}{2}

\noindent\textbf{第62題 不及底勾邊股條二}

\wen{乙從城外西南坤隅南行三百六十步，甲出北門東行二百步見之。問城徑。}
\da{城徑二百四十步。}
\shu{二行相乘即半徑冪。}
\shi{此以底勾大差股立法測望也。乙從坤隅南行大差股也，甲東行底勾也。底勾為城北半勾，大差股為城西南虛股。}

\switchcolumn

\noindent\textbf{Problem 62: Edge-Leg and Base-Leg Deficiency, 2}

\wenE{B walks south 360 paces from the Kun corner (southwest) outside the city; A exits the north gate and walks east 200 paces and sees B. Find the city's diameter.}
\daE{The city diameter is 240 paces.}
\shuE{Multiplying the two walks gives the square of the radius.}
\shiE{This establishes the method using base-leg and great-difference altitude for observation. B's southward walk from the Kun corner is the great-difference altitude; A's eastward walk is the base-leg. The base-leg is the northern half-base of the city; the great-difference altitude is the void-altitude southwest of the city.}

\end{paracol}

\bigskip

\begin{paracol}{2}

\noindent\textbf{第63題 不及底勾邊股條三}

\wen{甲出北門東行底勾也。底勾為城北半勾，明股為城南餘股。}
\da{城徑二百四十步。}
\shu{以底勾明股之數，約之得圓徑。}
\shi{此底勾明股立法也。以底勾明股之數，測望得圓徑。}

\switchcolumn

\noindent\textbf{Problem 63: Edge-Leg and Base-Leg Deficiency, 3}

\wenE{A exits the north gate and walks east---this is the base-leg. The base-leg is the northern half-base of the city; the bright-altitude is the remaining altitude south of the city.}
\daE{The city diameter is 240 paces.}
\shuE{Using the base-leg and bright-altitude values and reducing them gives the circle's diameter.}
\shiE{This establishes the method using base-leg and bright-altitude. Using these values, observation yields the circle's diameter.}

\end{paracol}

\clearpage

%% ========================================================================
%% Chapter 7: Shesheng buta chengfa
%% ========================================================================
\subsection*{測望不迭乘法 / Non-Iterative Observation Multiplication}
\addcontentsline{toc}{subsection}{測望不迭乘法 / Non-Iterative Observation Multiplication}

\begin{paracol}{2}

\noindent\textbf{總論：} 測望不迭乘法者，以測望之法，不迭次相乘，求圓城之徑也。

\switchcolumn

\noindent\textbf{General Discussion:} The non-iterative observation multiplication method uses the observation method without successive multiplications, to find the diameter of the circular city.

\end{paracol}

\bigskip

\begin{paracol}{2}

\noindent\textbf{第64題 測望不迭乘法一}

\wen{此法分別從方從廉明白，故重錄附之。}
\da{得半徑一百二十步。}
\shu{以測望之法，分別從方從廉，不迭次相乘，開平方得圓徑。}
\shi{此測望不迭乘法也。以測望之法，不迭次相乘，開平方得圓徑。}

\switchcolumn

\noindent\textbf{Problem 64: Non-Iterative Observation Multiplication, 1}

\wenE{This method clearly distinguishes the associate-square and associate-edge, so it is recorded again here for reference.}
\daE{The radius obtained is 120 paces.}
\shuE{Using the observation method, clearly distinguishing the associate-square and associate-edge, without successive multiplications, opening the square gives the circle's diameter.}
\shiE{This is the non-iterative observation multiplication method. Using the observation method without successive multiplications and opening the square gives the circle's diameter.}

\end{paracol}

\clearpage

%% ========================================================================
%% Reference Table for Volume 3
%% ========================================================================
\subsection*{卷三名數表 / Volume 3 Reference Table}
\addcontentsline{toc}{subsection}{卷三名數表 / Volume 3 Reference Table}

\begin{center}
\textbf{Chinese Terms \& Numerical Values}
\end{center}

\begin{center}
\begin{longtable}{|l|l|l|}
\hline
\textbf{Name} & \textbf{Value} & \textbf{Note} \\
\hline
\endfirsthead
\hline
\textbf{Name} & \textbf{Value} & \textbf{Note} \\
\hline
\endhead
上法 (shang fa) & First quotient squared & Upper position in root extraction \\
\hline
從方 (cong fang) & Next quotient $\times$ edge-method & Associate-square value \\
\hline
廉法 (lian fa) & Double corner-method & Edge-method value \\
\hline
隅法 (yu fa) & Next quotient squared & Corner-method value \\
\hline
下法 (xia fa) & Sum of square, edge, corner & Lower-method value \\
\hline
實 (shi) & Accumulated number & Dividend value \\
\hline
益實 (yi shi) & Added number & Augmenting dividend \\
\hline
減實 (jian shi) & Subtracted number & Subtracting dividend \\
\hline
負積 (fu ji) & Negative accumulation & Negative accumulated value \\
\hline
帶從 (dai cong) & Carrying accompanying-square & Accompanying-associate value \\
\hline
\end{longtable}
\end{center}

\clearpage

%% ========================================================================
%% END MATTER
%% ========================================================================
\section*{跋 / Postscript}
\addcontentsline{toc}{section}{跋 / Postscript}
\markboth{跋 / Postscript}{跋 / Postscript}

\begin{paracol}{2}

測圓海鏡分類釋術十卷，元李冶撰，明顧應祥分類釋術。其書以勾股測圓之術，分類編次，以便學者。卷一至卷三論通勾股、邊勾股、底股弦、皇極勾股、太陰勾股、明勾股、重勾股諸法求容圓，及兩股求容圓、兩弦求容圓、帶從開平方、減從翻法開平方、三乘方法、益積減積法、從方從廉隅法諸術。

\switchcolumn

The \textit{Ceyuan Haijing Fenlei Shishu}, in ten volumes, was composed by Li Ye of the Yuan Dynasty and classified with explanations by Gu Yingxiang of the Ming Dynasty. This book organizes the methods of right-triangle circle measurement by category and order, to facilitate students. Volumes 1--3 treat the methods of finding inscribed circles via universal, edge, base, supreme-pole, grand-yin, bright, and double right triangles, as well as the methods of two altitudes, two hypotenuses, accompanying-associate root extraction, subtractive-associate inverted root extraction, triple multiplication, additive and subtractive accumulation, and associate-square with associate-edge and corner methods.

\switchcolumn

其術之大要，以勾股相乘倍之為實，勾股和為法，除之得圓徑。又以弦和較即容圓徑。其開方之法，有帶從、減從翻法、益積、減積諸術，皆天元術之支流也。

\switchcolumn

The essentials of these methods are: multiply the base by the altitude and double it as the dividend; the sum of base and altitude is the divisor; dividing gives the circle's diameter. Also, the hypotenuse-sum difference is the inscribed circle's diameter. The methods of root extraction include accompanying-associate, subtractive-associate inverted, additive accumulation, and subtractive accumulation---all branches of the celestial element method (\textit{tian yuan shu}).

\switchcolumn

顧應祥序謂：李冶原書每條下細草，俱徑立天元一，反覆合之，而無下手之術。使後學之士茫然無門路可入。故去其細草，立一算術，又以其所立通勾邊股之屬各以類分之。此誠後學之津梁也。

\switchcolumn

As Gu Yingxiang states in his preface: In Li Ye's original work, beneath each entry the detailed working all proceeded directly by establishing the celestial element unity, folding and unfolding it repeatedly, without showing the method by which to begin. This left later students lost, with no door or path by which to enter. Therefore he removed the detailed working and established an explicit arithmetic procedure, also classifying the various quantities such as universal base and edge altitude by category. This is truly a bridge for later students.

\switchcolumn

\bigskip
\begin{flushright}
測圓海鏡分類釋術卷三終
\end{flushright}

\switchcolumn

\bigskip
\begin{flushright}
End of Volume 3 of the\\
\textit{Ceyuan Haijing Fenlei Shishu}
\end{flushright}

\end{paracol}

\clearpage

%% ========================================================================
%% COLOPHON
%% ========================================================================
\thispagestyle{empty}
\begin{center}
\vspace*{2.5cm}
{\Large 測圓海鏡分類釋術}

\vspace{0.3cm}
{\normalsize\textit{Ceyuan Haijing Fenlei Shishu}}

\vspace{0.8cm}
{\large 卷一至卷三終}

\vspace{0.3cm}
{\normalsize\textit{End of Volumes 1--3}}

\vspace{1.5cm}
\shortline

\vspace{0.8cm}
{\normalsize
\begin{tabular}{rl}
總校官候補中書 & 臣吳紹潔 \\
\textit{Chief Proofreader:} & \textit{Wu Shaojie} \\
\[0.3em]
校對官中官正 & 臣郭長發 \\
\textit{Proofreader:} & \textit{Guo Changfa} \\
\[0.3em]
謄錄監生 & 臣丁湘錦 \\
\textit{Copyist:} & \textit{Ding Xiangjin} \\
\end{tabular}
}

\vspace{2cm}
{\small\textcolor{notecolor}{%
此為欽定四庫全書薈要本\\
子部天文算法類\\
清乾隆年間鈔本
}}

\vspace{0.3cm}
{\small\textcolor{notecolor}{%
\textit{Imperial edition from the Siku Quanshu Huiyao}\\
\textit{Mathematical and Astronomical Works Division}\\
\textit{Manuscript copy from the Qianlong period (1736--1795)}
}}

\end{center}

\clearpage

%% ========================================================================
%% MODERN note
%% ========================================================================
\section*{Modern Note / 現代編輯說明}
\addcontentsline{toc}{section}{Modern Note / 現代編輯說明}
\markboth{Modern Note / 現代編輯說明}{Modern Note / 現代編輯說明}

\begin{paracol}{2}

本版為《測圓海鏡分類釋術》卷一至卷三雙語對照本。原書為元代李冶所撰《測圓海鏡》，明代顧應祥分類釋術，收入《四庫全書薈要》。

\switchcolumn

This edition presents a bilingual English--Chinese version of Volumes 1--3 of the \textit{Ceyuan Haijing Fenlei Shishu} (Classified Methods of the Sea Mirror of Circle Measurements). The original work was Li Ye's \textit{Ceyuan Haijing}, rearranged and explained by Gu Yingxiang during the Ming dynasty, included in the \textit{Siku Quanshu Huiyao}.

\switchcolumn

\subsection*{關於原著 / About the Original Work}

原《測圓海鏡》為元李冶（1192--1279）於1248年所著。此書為中國古代數學之瑰寶，以勾股容圓為核心，凡170題，運用天元術（古代代數方法）解決各種幾何問題。

\switchcolumn

\subsection*{About the Original Work}

The original \textit{Ceyuan Haijing} was composed by Li Ye (1192--1279) in 1248 during the Yuan dynasty. It is one of the masterpieces of classical Chinese mathematics, centering on right-triangle inscribed-circle problems with 170 exercises, employing the \textit{tian yuan shu} (celestial element method---an ancient algebraic technique) to solve various geometric problems.

\switchcolumn

\subsection*{顧應祥之分類釋術 / Gu Yingxiang's Rearrangement}

明代顧應祥（1483--1565）認為李冶原書雖精妙，但其細草皆用天元術，使後學難以入門。故顧氏去其細草，立一算術，又以數學方法分類重編，加釋術以明其理。

\switchcolumn

\subsection*{Gu Yingxiang's Rearrangement}

Gu Yingxiang (1483--1565) of the Ming Dynasty found Li Ye's original work, though brilliant, difficult for students because its detailed working all employed the \textit{tian yuan shu}. Therefore, Gu removed the detailed working, established explicit arithmetic procedures, reorganized the problems by mathematical method, and added explanatory notes to clarify the reasoning.

\switchcolumn

\subsection*{卷一至卷三結構 / Structure of Volumes 1--3}

\begin{itemize}
\item \textbf{卷一}論通勾股、邊勾股、底股弦、皇極勾股、太陰勾股、明勾股、重勾股求容圓之法。
\item \textbf{卷二}論兩股求容圓、兩弦求容圓、長短二勾求城徑、減從翻法開平方、勾股相乘倍之開平方、弦和較求圓徑之法。
\item \textbf{卷三}論帶從開平方、減從翻法開平方、三乘方、益積減積法、從方從廉隅法、不及底勾邊股條、測望不迭乘法之法。
\end{itemize}

\switchcolumn

\subsection*{Structure of Volumes 1--3}

\begin{itemize}
\item \textbf{Volume 1} treats inscribed-circle methods via universal, edge, base, supreme-pole, grand-yin, bright, and double right triangles.
\item \textbf{Volume 2} covers methods using two altitudes, two hypotenuses, long and short bases, subtractive-associate inverted root extraction, doubled-product root extraction, and hypotenuse-sum difference.
\item \textbf{Volume 3} presents accompanying-associate root extraction, subtractive-associate inverted root extraction, triple multiplication, additive and subtractive accumulation, associate-square with edge and corner methods, edge-leg and base-leg deficiency, and non-iterative observation multiplication.
\end{itemize}

\switchcolumn

\subsection*{數學意義 / Mathematical Significance}

全書諸題皆圍繞一核心幾何構造：一圓城，自各觀測點量得之數據各異，然所求之城徑皆為\textbf{二百四十步}。此非巧合，乃展示代數結構之統一性也。

\switchcolumn

\subsection*{Mathematical Significance}

All problems in the book revolve around a single geometric configuration: a circular city with varying measurements from different observation points, yet the city diameter sought is invariably \textbf{240 paces}. This is no coincidence but rather demonstrates the algebraic unity underlying all the problems.

\switchcolumn

其方法展現中國古代高度發達之開方術與多項式代數傳統，遠早於歐洲類似發展數百年。顧應祥之分類系統揭示了李冶170題背後之結構統一性，按數學方法而非表面幾何形態加以組織。

\switchcolumn

These methods exemplify the sophisticated Chinese tradition of root extraction and polynomial algebra that predated comparable European developments by centuries. Gu Yingxiang's classification system reveals the underlying structural unity of Li Ye's 170 problems, organizing them by mathematical technique rather than surface geometric appearance.

\end{paracol}

\vspace{1cm}
\begin{flushright}
\textit{Bilingual edition typeset with XeLaTeX using Noto CJK fonts}\\
\textit{Modern edition prepared from the Siku Quanshu Huiyao manuscript}\\
\textit{雙語版使用 XeLaTeX 及 Noto CJK 字體排版}\\
\textit{現代版據《四庫全書薈要》鈔本製作}
\end{flushright}

\end{document}
